\chapter[Lie-Algebra Links as Latent Fields]{Lie-Algebra Links as Latent Fields}
\label{ch:links}

This chapter asks what becomes of the independent links of \Cref{def:geo-graph-links} when they are
inferred jointly with the agents' states instead of being held as structural data or moved as
parameters. Written in the exponential chart of a reference point, every twist is the exponential of a
generator in the Lie algebra, on which passive re-expressions do not act
(\Cref{def:link-exponential-chart,prop:link-generator-gauge}). A Gaussian base law on the generators
multiplied by the exponential of minus a nonnegative loop action is a normalized prior for every
structure group whose Lie algebra carries an inner product, compact or not
(\Cref{prop:link-generator-prior}); on the general linear group the loop action measures the reference
holonomy with a fiber metric, because no continuous class function separates unipotent holonomy from
the identity (\Cref{cor:link-no-invariant-defect,prop:link-loop-defect}). With the generators latent,
the collective bound of \Cref{prop:elbo-random-link-vfe} specializes to a joint that has unit mass on
every record hypergraph and to a free energy whose link sector is the relative entropy from the
generator prior (\Cref{prop:link-latent-normalization,prop:link-gaussian-recognition-vfe}). Its
stationarity condition is a link field equation in which the stress of the prior balances the expected
matter current transported by the adjoint derivative of the chart (\Cref{prop:link-field-equation});
for a comparison record the current has the outer-product form of a delta rule
(\Cref{cor:link-delta-rule}), and without the prior the link problem degenerates
(\Cref{cor:link-prior-necessary}). Held links are the zero-width limit of the latent model, and
statements proved at every held link transfer conditionally on the generators
(\Cref{thm:link-held-limit,prop:link-conditional-transfer}). Matter attached by record extension
induces no link action, while matter normalized jointly with the generators induces minus the logarithm
of its normalizer (\Cref{prop:link-induced-action}). Recognition of the generators is a first-order
natural-gradient history whose explicit Euler step is the Gaussian Bayesian learning rule
(\Cref{prop:link-recognition-flow}); across learning episodes a Gibbs prior on generator histories gives
the most probable history an inertia whose mass is the precision of the generator increment
(\Cref{prop:link-history-prior,prop:link-map-inertia}); and near zero generators the stationarity
condition is a linear equation whose positive operator collects the anchor, the loops, the episodes,
and the matter (\Cref{prop:link-linearized-equation}). These results are used for the coarse links and
the holonomy-conditioned parents of \Cref{ch:network-state,ch:full-graph-meta-agent-vfe}, whose link
statements gain a version with latent generators, and the link translations of
\Cref{ch:finite-depth-formation} are their translation-group instance
(\Cref{prop:link-translation-instance}); Supplement G realizes them for Gaussian matter. Longer proofs
are deferred to \cref{app:deferred-proofs}.

The symbols of \cref{tab:link-notation} are declared for this chapter and its proof section, and those
that recur in \Cref{ch:network-state,ch:full-graph-meta-agent-vfe} are also listed in \cref{tab:notation};
$\mathbf B$, $\mathbf S$, $\mathbf T$, and $\mathbf H$ denote generic square matrices, $\mathbf b$ and
$\mathbf c$ generic vectors, and $\mathbf h$ a generic real function, each bound in the statement or proof
where it appears, and a ring accent marks an object of the generator sector written in the coordinates of
the reference point.

\begin{sciencetable}{Symbols of the latent-link construction.}
\label{tab:link-notation}
\begin{tabularx}{\textwidth}{>{\raggedright\arraybackslash}p{0.33\textwidth}>{\raggedright\arraybackslash}X}
\toprule
\textbf{Symbol} & \textbf{Meaning} \\
\midrule
\(\mathfrak g\), \(\exp\), \(\operatorname{Ad}_g\), \(\mathrm d\rho_x\) & Lie algebra of \(G\), its exponential map, the adjoint action, the differential of \(\rho_x\) at \(e_G\) \\
\(E_+\), \(\mathscr T=\mathfrak g^{E_+}\) & one orientation of each reversal pair; link space of one channel, \(\mathscr T\times\mathscr T\) for two independent channels \\
\(\mathring\omega^x_e\), \(\exp_\sigma\), \(\mathring\Theta^x_p\) & generator of the twist of \(e\); exponential chart at \(\sigma\); reference path product, the reference holonomy of a closed walk \\
\(\mathsf G^{\mathrm{fib}}\), \(\langle\cdot,\cdot\rangle_{\mathsf G}\), \(\lVert\cdot\rVert_{\mathsf G}\) & fiber metric at the reference point; trace form on \(\mathfrak{gl}(\mathbb V)\) and the inner product on \(\mathfrak g\) \\
\(\mathfrak Q^x_p\), \(\mathbb L\), \(\mathrm c_p\), \(\mathbf C\) & loop defect of a closed walk; family of closed walks and their coefficients; cycle map \\
\(\mathbb L_+\), \(\mathbf C_+\), \(z_p\), \(r_{\mathbb L}\), \(\kappa\) & loops with positive coefficient and their cycle map; net traversal vector of a loop; rank of these vectors; number of components of the multigraph \\
\(d\), \(Z\), \(\mathscr T_{\mathrm{ex}}\), \(\mathscr T_{\mathrm{har}}\) & vertex-difference map and cycle space of \(\R^{E_+}\); exact and harmonic generators \\
\(\nu^{\mathscr T}\), \(\mathscr S\), \(Z_{\mathscr S}\), \(\mathbb P^{\mathscr T}_\theta\) & base law, link action, normalizer, generator prior \\
\(\mathring\alpha^x_e\), \(\mathring\Sigma^x_e\); \(\mathring\mu\), \(\mathring\Lambda\) & anchor and prior covariance; mean and precision of a Gaussian recognition law \(R\) of the generators \\
\(\mathbf L(\mathbf B,\mathbf T)\) & Fr\'echet derivative of the matrix exponential \\
\(\mathbf H_o\), \(\mathbf U\), \(\mathbf V\), \(\mathbf J^x_e\) & matter energy, expected action plus matter energy, most-probable-link objective, link current \\
\(\mathbf M\), \(\mathbf M^\star\), \(\mathbf l\) & transport matrix of a comparison record, its rank-one minimizer, expected record energy \\
\(\mathbf F_Q\), \(\mathbf F\) & held-link free energy as a function of the generators, and its infimum over a class \\
\(\mathscr S^{\mathrm{ind}}\), \(\mathfrak Z\) & induced link action, joint normalizer of generators and undirected matter \\
\(\mathfrak i_\lambda\), \(\mathring\Xi\), \(\mathring\vartheta\), \(\mathring\varpi\), \(\mathring\Gamma\) & generator integrand, space of admissible precisions, natural and expectation parameters, log-partition function \\
\(\mathbf m\), \(\mathfrak U\), \(\mathbb P^{\mathrm{hist}}\), \(Z_{\mathrm{hist}}\), \(\mathring Z^{\leftarrow}_n\) & increment precision, kinetic term, generator history prior, its normalizer, backward functions \\
\(\mathscr S^{\mathrm{hist}}\), \(\mathfrak V_n\) & history action and episode potentials \\
\(\mathbf K\), \(\mathbf N_0\), \(\mathring\Phi\), \(\mathbf j\) & linearized link operator, current correction of the Hessian, matter operator, source \\
\bottomrule
\end{tabularx}
\end{sciencetable}

\paragraph{Scope.}
The chapter does not claim that latent links improve the evidence, the prediction of any record, or any
learning outcome relative to held links; that the generator prior is invariant under reference changes
that do not act by multiples of isometries of the fiber metric, or that the loop defect is independent
of the base vertex of a closed walk; that every local minimizer of the most-probable-link objective is
global; that the recognition flow of the generators converges, or that stationary generator histories
are unique outside the ball of \cref{prop:link-linearized-equation}; that the joint M-coordinate of the
anchors and the prior covariances exists when the action has loop terms; that frames are dynamical, so
that time-like links would carry variables; that the linearized link operator is the operator of a
physical field; or that the outer-product form of the link current identifies the model with a
fast-weight or delta-rule architecture. \status{NOT-CLAIMED} The interaction multigraph stays structural, and edge creation and
deletion is \cref{open:obs-edge-creation}; a finite-dimensional family of link laws that survives the
coarse step is \cref{open:link-coarse-closure}; and a slow-manifold reduction with the generator
recognition means in the slow block is part of \cref{open:obs-slow-manifold}.

\section{The exponential chart of the links}
\label{sec:link-chart}

This section writes the independent links of \Cref{def:geo-graph-links} in the exponential chart of
the reference point $\sigma$ of \Cref{prop:geo-edge-twist-invariance}. The twist of an edge copy is
unchanged by every passive re-expression, so a generator of the twist carries no vertex gauge: the
frames $U_i$ carry the passive action, while a change of reference point or an active automorphism
moves all generators by one adjoint map (\Cref{prop:link-generator-gauge}). The chart covers exactly
the link families whose twists are exponentials. For $G=\GL(K,\R)$ this excludes some twists of
positive determinant, while on a closed walk that traverses two reversal pairs once each the
reference holonomies of chart links still reach every element of $\GL^+(K)$
(\Cref{prop:link-chart-image}).

\lemmaheading{Exponential maps and homomorphisms}{lem:link-exp-naturality}
Let $\boldsymbol\varphi:G\to\mathbb Y$ be a smooth homomorphism of Lie groups with differential
$\mathrm d\boldsymbol\varphi:\mathfrak g\to\operatorname{Lie}(\mathbb Y)$ at $e_G$. Then
$\boldsymbol\varphi(\exp\mathring\omega)=\exp(\mathrm d\boldsymbol\varphi\mathring\omega)$ for every
$\mathring\omega\in\mathfrak g$. In particular, let $\operatorname{Ad}_g$ be the differential at $e_G$
of the conjugation $g'\mapsto gg'g^{-1}$, and let $\rho_x:G\to\GL(\mathbb V)$ be a smooth homomorphism
into the general linear group of a finite-dimensional real vector space $\mathbb V$, with differential
$\mathrm d\rho_x$ at $e_G$. Then for all $g\in G$ and $\mathring\omega\in\mathfrak g$, with the matrix
exponential as the exponential map of $\GL(\mathbb V)$,
\begin{equation}
g\exp(\mathring\omega)g^{-1}=\exp(\operatorname{Ad}_g\mathring\omega),
\qquad
\rho_x(\exp\mathring\omega)=\exp(\mathrm d\rho_x\mathring\omega),
\qquad
\mathrm d\rho_x(\operatorname{Ad}_g\mathring\omega)=\rho_x(g)\mathrm d\rho_x(\mathring\omega)\rho_x(g)^{-1}.
\label{eq:link-exp-naturality}
\end{equation}
\status{ESTABLISHED}

\noindent\emph{Proof.}
The exponential map sends $\mathring\omega$ to the value at time one of the integral curve through
the identity of the left-invariant vector field whose value at the identity is $\mathring\omega$, and
$t\mapsto\exp(t\mathring\omega)$ is a smooth homomorphism from $(\R,+)$. Hence
$\boldsymbol\gamma(t)=\boldsymbol\varphi(\exp(t\mathring\omega))$ is a smooth homomorphism from $(\R,+)$ into $\mathbb Y$.
Differentiating $\boldsymbol\gamma(t+t')=\boldsymbol\gamma(t)\boldsymbol\gamma(t')$ in $t'$ at $t'=0$ shows that $\boldsymbol\gamma'(t)$ is the image
of $\boldsymbol\gamma'(0)$ under the differential at the identity of the left multiplication by $\boldsymbol\gamma(t)$, so
$\boldsymbol\gamma$ is the integral curve through the identity of the left-invariant vector field with value
$\boldsymbol\gamma'(0)=\mathrm d\boldsymbol\varphi\mathring\omega$, and uniqueness of integral curves gives
$\boldsymbol\gamma(t)=\exp(t\mathrm d\boldsymbol\varphi\mathring\omega)$; the first claim is the case $t=1$. Conjugation by
$g$ is a smooth homomorphism with differential $\operatorname{Ad}_g$, which gives the first identity of
\cref{eq:link-exp-naturality}. On $\GL(\mathbb V)$ the left-invariant vector field with value
$\mathbf H$ at the identity is $y\mapsto y\mathbf H$, and the matrix exponential satisfies
$\frac{\mathrm d}{\mathrm dt}e^{t\mathbf H}=e^{t\mathbf H}\mathbf H$ with $e^{0}=I_{\mathbb V}$, so it is the
exponential map of $\GL(\mathbb V)$; this gives the second identity. Applying $\rho_x$ to the first
identity for $t\mathring\omega$ and using the second gives
$\rho_x(g)\exp(t\mathrm d\rho_x\mathring\omega)\rho_x(g)^{-1}=\exp(t\mathrm d\rho_x(\operatorname{Ad}_g\mathring\omega))$,
whose derivative at $t=0$ is the third identity. \hfill$\square$

\definitionheading{Exponential chart of the links}{def:link-exponential-chart}
Let $\mathfrak g$ be the Lie algebra of $G$ and $\exp:\mathfrak g\to G$ its exponential map, let
$E_+\subseteq E$ contain exactly one edge copy of each reversal pair of $\mathfrak I$, and let the
reference point $\sigma$ and the frames $u_i=\sigma\cdot U_i^{-1}$ be those of
\Cref{prop:geo-edge-twist-invariance}. For a channel $x$, a generator family is a point
$\mathring\omega^x=(\mathring\omega^x_e)_{e\in E_+}$ of the link space $\mathscr T=\mathfrak g^{E_+}$,
extended to the reversed copies by the first line below. Its twists and links are those of the second
line, and for a walk $p=(e_0,\ldots,e_{r-1})$ with $e_a:i_{a+1}\to i_a$, as in
\Cref{def:geo-graph-links}, its reference path product is that of the third line, called the reference
holonomy of $p$ when $p$ is closed:
\begin{gather}
\mathring\omega^x_{\bar e}=-\mathring\omega^x_e,
\label{eq:link-generator-reversal}\\
\mathring\Theta^x_e=\exp(\mathring\omega^x_e),
\qquad
\Theta^x_e=U_i\exp(\mathring\omega^x_e)U_j^{-1}\quad(e:j\to i),
\label{eq:link-chart}\\
\mathring\Theta^x_p=\mathring\Theta^x_{e_0}\mathring\Theta^x_{e_1}\cdots\mathring\Theta^x_{e_{r-1}}
=\exp(\mathring\omega^x_{e_0})\exp(\mathring\omega^x_{e_1})\cdots\exp(\mathring\omega^x_{e_{r-1}}).
\label{eq:link-reference-holonomy}
\end{gather}
The links satisfy \cref{eq:geo-regime-two-links} because $\exp(-\mathring\omega)=\exp(\mathring\omega)^{-1}$,
and $\mathring\Theta^x_e$ is their twist in the sense of \cref{eq:geo-edge-twist}. The exponential chart
at $\sigma$ is the continuous map $\exp_\sigma:\mathscr T\to G^{E_+}$,
$\mathring\omega^x\mapsto(U_i\exp(\mathring\omega^x_e)U_j^{-1})_{e\in E_+}$. The link space is a real
vector space of dimension $\lvert E_+\rvert\dim G$, standard Borel with its Borel sets, and the chart
covers exactly the link families whose twists lie in $\exp(\mathfrak g)$ (\Cref{prop:link-chart-image}).
When $\rho_x$ is a smooth homomorphism $G\to\GL(\mathbb V)$ on a real vector sample fiber
$\mathbb V=\R^{\dim\mathbb V}$, the represented generators are
$\mathrm d\rho_x(\mathring\omega^x_e)\in\mathfrak{gl}(\mathbb V)$, and
$\rho_x(\mathring\Theta^x_e)=\exp(\mathrm d\rho_x\mathring\omega^x_e)$ by \Cref{lem:link-exp-naturality}.
The two channels carry independent generator families, with link space $\mathscr T\times\mathscr T$;
under the common-link specialization \cref{eq:geo-common-link} one generator family
$\mathring\omega=\mathring\omega^b=\mathring\omega^m$ serves both channels and the index $x$ is dropped.

\propositionheading{Gauge action on generators}{prop:link-generator-gauge}
In \Cref{def:link-exponential-chart}, let the generator family $\mathring\omega^x$ have the links
$\Theta^x=\exp_\sigma(\mathring\omega^x)$. (a) After a passive re-expression $(g_i)_{i\in V}$ the links
$g_i\Theta^x_eg_j^{-1}$ are the links of the same $\mathring\omega^x$ relative to $\sigma$ and the
re-expressed frames $g_iU_i$, so every generator is unchanged. (b) Relative to the reference point
$\sigma'=\sigma\cdot g_0$ the frames are $U_ig_0$, and the links $\Theta^x$ are those of the generator
family $\operatorname{Ad}_{g_0^{-1}}\mathring\omega^x=(\operatorname{Ad}_{g_0^{-1}}\mathring\omega^x_e)_{e\in E_+}$,
so that $\exp_{\sigma'}\circ\operatorname{Ad}_{g_0^{-1}}=\exp_\sigma$; every reference path product
becomes $g_0^{-1}\mathring\Theta^x_pg_0$. (c) An active $F$ with $F(\sigma)=\sigma\cdot\kappa$, as in
\Cref{prop:geo-edge-twist-invariance}, carries $\Theta^x$ to the links of
$\operatorname{Ad}_\kappa\mathring\omega^x$ relative to $\sigma$ and the same frames, and every
reference path product to $\kappa\mathring\Theta^x_p\kappa^{-1}$. (d) The holonomy of a closed walk
$p$ based at $i_0$ and, for a smooth homomorphism $\rho_x:G\to\GL(\mathbb V)$, its represented
reference holonomy are
\begin{equation}
\mathrm{Hol}^x(p)=U_{i_0}\mathring\Theta^x_pU_{i_0}^{-1},
\qquad
\rho_x(\mathring\Theta^x_p)=\exp(\mathrm d\rho_x\mathring\omega^x_{e_0})\exp(\mathrm d\rho_x\mathring\omega^x_{e_1})\cdots\exp(\mathrm d\rho_x\mathring\omega^x_{e_{r-1}}).
\label{eq:link-holonomy-chart}
\end{equation}
Each of (a) to (c) moves all generators by one linear map of $\mathfrak g$, namely the identity,
$\operatorname{Ad}_{g_0^{-1}}$, or $\operatorname{Ad}_\kappa$, and the passive vertex group acts on the
frames alone. \status{ESTABLISHED}

\noindent\emph{Proof.}
(a) The vertex law \cref{eq:geo-regime-two-gauge-law} and $U_i'=g_iU_i$ from
\cref{eq:geo-regime-one-gauge-law} give
\begin{equation*}
g_i\Theta^x_eg_j^{-1}=(g_iU_i)\exp(\mathring\omega^x_e)(g_jU_j)^{-1}.
\end{equation*}
(b) From
$\sigma=\sigma'\cdot g_0^{-1}$ we get $u_i=\sigma'\cdot(U_ig_0)^{-1}$, and
\Cref{lem:link-exp-naturality} gives
$(U_ig_0)\exp(\operatorname{Ad}_{g_0^{-1}}\mathring\omega^x_e)(U_jg_0)^{-1}=U_ig_0g_0^{-1}\exp(\mathring\omega^x_e)g_0g_0^{-1}U_j^{-1}=\Theta^x_e$;
the same identity applied to each factor of \cref{eq:link-reference-holonomy} gives the path products.
(c) By \Cref{prop:geo-edge-twist-invariance} the image twists are
$\kappa\exp(\mathring\omega^x_e)\kappa^{-1}=\exp(\operatorname{Ad}_\kappa\mathring\omega^x_e)$, and a twist
determines its link given the frames through $\Theta^x_e=U_i\mathring\Theta^x_eU_j^{-1}$. (d) The first
identity is \cref{eq:geo-edge-twist-holonomy}, and the second follows from the homomorphism property of
$\rho_x$ and \Cref{lem:link-exp-naturality}. \hfill$\square$

\propositionheading{Real logarithms}{prop:link-real-logarithm}
A real square matrix is the exponential of a real square matrix exactly when it is nonsingular and
each Jordan block belonging to a negative eigenvalue occurs an even number of times
\citep[Theorem~1]{culver1966existence}. \status{STANDARD}

\propositionheading{Image of the exponential chart}{prop:link-chart-image}
In \Cref{def:link-exponential-chart} the following hold. (a) The image of $\exp_\sigma$ is the set of
link families of channel $x$ whose twists all lie in $\exp(\mathfrak g)$; it does not depend on
$\sigma$, and passive re-expressions and active automorphisms map it onto itself. Link families with a
twist outside $\exp(\mathfrak g)$ have no generator family. (b) For $G=\GL(K,\R)$ or $G=\GL^+(K)$ the
set $\exp(\mathfrak{gl}(K,\R))$ is the set of matrices described in \Cref{prop:link-real-logarithm}
and lies in $\GL^+(K)$; for $K\geq2$ it is a proper subset of $\GL^+(K)$, since
$\operatorname{diag}(-1,-2,1,\ldots,1)$ is not the exponential of a real matrix, so the chart misses
every link family with such a twist, and for $G=\GL(K,\R)$ it also misses every link family with a
twist of negative determinant. (c) Every element of $\GL^+(K)$ equals $\exp(\mathbf S)\exp(\mathbf T)$
with $\mathbf S$ symmetric and $\mathbf T$ skew-symmetric, and every element of $SO(K)$ equals
$\exp(\mathbf T)$ with $\mathbf T$ skew-symmetric; hence $\GL^+(K)$ is connected and generated by
$\exp(\mathfrak{gl}(K,\R))$, and for $G=SO(K)$ the chart covers every link family. (d) For every
connected Lie group $G$, $\exp(\mathfrak g)$ generates $G$. (e) Let $G$ be $\GL(K,\R)$ or $\GL^+(K)$,
and let the closed walk $p$ traverse two distinct reversal pairs exactly once each. Then every element
of $\GL^+(K)$ is the reference holonomy $\mathring\Theta_p$ of some generator family, so for $K\geq2$
the reference holonomies of chart links leave $\exp(\mathfrak{gl}(K,\R))$. On a triangle walk
$(e_0,e_1,e_2)$ with every $e_a\in E_+$ and $K=2$, the generators $\operatorname{diag}(0,\log2)$,
$\pi(E_{21}-E_{12})$, and $0$ give $\mathring\Theta_p=\operatorname{diag}(-1,-2)$. \status{ESTABLISHED}

The proof is in \cref{sec:proofs-links}.

\section{Curvature on a noncompact structure group}
\label{sec:link-curvature}

A loop penalty should vanish exactly on the closed walks with trivial represented holonomy. For compact
$G$ and unitary $\rho_x$ the class function of \Cref{rem:geo-loop-penalties} does so; on the general
linear group no continuous class function can (\Cref{cor:link-no-invariant-defect}), because
conjugation contracts every unipotent element to the identity
(\Cref{prop:link-class-function-blindness}). The loop defect defined below measures the reference
holonomy with a fiber metric placed at the reference point. It separates every holonomy from the
identity; with the metric held as a matrix at the reference point, it is invariant under the reference
changes that act by multiples of isometries of that metric, and for the general linear group under no
others (\Cref{prop:link-loop-defect}).

\propositionheading{Continuous class functions are blind to unipotent holonomy}{prop:link-class-function-blindness}
Let $K\geq2$, let $\mathbb H$ be $\GL(K,\R)$ or $\GL^+(K)$, and let $\mathbf h:\mathbb H\to\R$ be
continuous and constant on the conjugacy classes of $\mathbb H$. Then $\mathbf h(I_K+\mathbf T)=\mathbf h(I_K)$
for every nilpotent $\mathbf T$, that is, at every unipotent matrix $I_K+\mathbf T$, which lies in
$\GL^+(K)$ because $\det(I_K+\mathbf T)=1$. \status{ESTABLISHED}

\noindent\emph{Proof.}
Since $\mathbf T^K=0$, the kernels
$\ker\mathbf T\subseteq\ker\mathbf T^2\subseteq\cdots\subseteq\ker\mathbf T^K=\R^K$ increase, and
$\mathbf T$ maps $\ker\mathbf T^{k+1}$ into $\ker\mathbf T^k$. Listing a basis of $\ker\mathbf T$,
extending it to a basis of $\ker\mathbf T^2$, and so on, gives the columns of an invertible $\mathbf B$
with $\mathbf B^{-1}\mathbf T\mathbf B=\mathbf T_0$ strictly upper triangular. Replacing the first column
of $\mathbf B$ by its negative conjugates $\mathbf T_0$ by a diagonal matrix, which keeps it strictly
upper triangular, so $\det\mathbf B>0$ may be assumed. For $t>0$ let
$\mathbf S_t=\operatorname{diag}(1,t,\ldots,t^{K-1})$; then $\mathbf B\mathbf S_t\in\GL^+(K)\subseteq\mathbb H$
and $(\mathbf S_t^{-1}\mathbf T_0\mathbf S_t)_{kl}=t^{l-k}(\mathbf T_0)_{kl}$, which tends to zero as
$t\to0$ because $(\mathbf T_0)_{kl}=0$ unless $l>k$. Hence
$\mathbf h(I_K+\mathbf T)=\mathbf h\bigl((\mathbf B\mathbf S_t)^{-1}(I_K+\mathbf T)\mathbf B\mathbf S_t\bigr)
=\mathbf h(I_K+\mathbf S_t^{-1}\mathbf T_0\mathbf S_t)$ for every $t>0$, and continuity gives
$\mathbf h(I_K+\mathbf T)=\mathbf h(I_K)$. \hfill$\square$

\corollaryheading{No invariant loop defect on the general linear group}{cor:link-no-invariant-defect}
Let $K\geq2$ and $\mathbb H$ be as in \Cref{prop:link-class-function-blindness}. No continuous function
on $\mathbb H$ that is constant on its conjugacy classes vanishes exactly at $I_K$: every such function
that vanishes at $I_K$ vanishes at the shear $I_K+E_{12}\neq I_K$, and so does every function of the
eigenvalues, counted with multiplicity, that vanishes at $I_K$. Consequently, for $G=\GL(K,\R)$ or
$G=\GL^+(K)$ and a closed walk $p$ that traverses some reversal pair exactly once, a loop penalty that
is a continuous function of $\mathrm{Hol}^x(p)$ on $G$ and is unchanged by every passive re-expression
of every link family takes at every unipotent holonomy the value it takes at the trivial holonomy, so it
vanishes at every unipotent holonomy whenever it vanishes at the trivial one. In particular no norm on
$\mathfrak{gl}(K,\R)$ is invariant under conjugation by $\mathbb H$, which for $\mathbb H=\GL(K,\R)$ is
the observation of \Cref{rem:geo-loop-penalties}; the corollary is stronger, since it excludes every
continuous conjugation-invariant penalty and not only those of the form $g'\mapsto\lVert g'-I_K\rVert$.
\status{ESTABLISHED}

\noindent\emph{Proof.}
The first claim is \Cref{prop:link-class-function-blindness} at $\mathbf T=E_{12}$, which is
nilpotent because $E_{12}^2=0$; the shear is upper triangular with unit diagonal, so its eigenvalues are
those of $I_K$. If $p$ traverses the pair of $e$ exactly once, the link family with $\Theta^x_e$ equal
to a given $g'\in G$ or to $g'^{-1}$, according to the direction of the traversal, and with all other
links equal to $e_G$ has $\mathrm{Hol}^x(p)=g'$, and the passive re-expression with $g_{i_0}=g$ replaces
it by $gg'g^{-1}$ by \cref{eq:geo-edge-twist-holonomy}; the penalty is therefore a continuous class
function on $G$. For a norm $\lVert\cdot\rVert_\star$ invariant under conjugation by $\mathbb H$, the
function $g'\mapsto\lVert g'-I_K\rVert_\star$ is continuous, constant on conjugacy classes because
$gg'g^{-1}-I_K=g(g'-I_K)g^{-1}$, and vanishes exactly at $I_K$. \hfill$\square$

\definitionheading{Loop defect and cycle map}{def:link-loop-defect}
In \Cref{def:link-exponential-chart}, let $\rho_x:G\to\GL(\mathbb V)$ be a smooth homomorphism on a real
vector sample fiber $\mathbb V=\R^{\dim\mathbb V}$, and let $\mathsf G^{\mathrm{fib}}$ be a symmetric
positive definite matrix at the reference point: the matrix, in the coordinates of $\sigma$, of an inner
product on the fiber at $c_*$ of the vector bundle associated with $\rho_x$. It plays the role of the
fiber metric of \Cref{rem:geo-loop-penalties}, placed at $\sigma$ instead of at a vertex. On
$\mathfrak{gl}(\mathbb V)$ put
\begin{equation}
\langle\mathbf H,\mathbf H'\rangle_{\mathsf G}=\Tr\bigl((\mathsf G^{\mathrm{fib}})^{-1}\mathbf H^\top\mathsf G^{\mathrm{fib}}\mathbf H'\bigr),
\qquad
\lVert\mathbf H\rVert_{\mathsf G}=\langle\mathbf H,\mathbf H\rangle_{\mathsf G}^{1/2},
\label{eq:link-trace-form}
\end{equation}
and, when $\mathrm d\rho_x$ is injective, write
$\langle\mathring\omega,\mathring\omega'\rangle_{\mathsf G}=\langle\mathrm d\rho_x\mathring\omega,\mathrm d\rho_x\mathring\omega'\rangle_{\mathsf G}$
for the pulled-back inner product on $\mathfrak g$. The loop defect of a closed walk $p$ based at $i_0$
is the first quantity below, with $\mathring\Theta^x_p$ the reference holonomy
\cref{eq:link-reference-holonomy}. For a finite family $\mathbb L$ of closed walks, the cycle map
$\mathbf C:\mathscr T\to\mathfrak g^{\mathbb L}$ is the second, with $p=(e_0,\ldots,e_{r-1})$ and the
reversal convention \cref{eq:link-generator-reversal}:
\begin{gather}
\mathfrak Q^x_p(\mathring\omega^x)=\tfrac12\bigl\lVert I_{\mathbb V}-\rho_x(\mathring\Theta^x_p)\bigr\rVert_{\mathsf G}^2,
\label{eq:link-loop-defect}\\
(\mathbf C\mathring\omega^x)_p=\mathring\omega^x_{e_0}+\mathring\omega^x_{e_1}+\cdots+\mathring\omega^x_{e_{r-1}}.
\label{eq:link-cycle-map}
\end{gather}
Each reversal pair thus enters $(\mathbf C\mathring\omega^x)_p$ with its generator on $E_+$ times the
number of traversals of $p$ along the orientation in $E_+$ minus the number against it, and $\mathbf C$
is linear.

\lemmaheading{The trace form determines the metric up to scale}{lem:link-trace-form-scale}
For symmetric positive definite matrices $\mathsf G'$ and $\mathsf G''$ on $\mathbb V$, write
$\lVert\cdot\rVert_{\mathsf G'}$ and $\lVert\cdot\rVert_{\mathsf G''}$ for \cref{eq:link-trace-form} with
$\mathsf G'$ and $\mathsf G''$ in place of $\mathsf G^{\mathrm{fib}}$. Then
$\lVert\mathbf H\rVert_{\mathsf G'}=\lVert\mathbf H\rVert_{\mathsf G''}$ for every rank-one $\mathbf H$ exactly
when $\mathsf G''$ is a positive multiple of $\mathsf G'$, and in that case the two norms agree on all of
$\mathfrak{gl}(\mathbb V)$. \status{ESTABLISHED}

\noindent\emph{Proof.}
For $\mathbf H=\mathbf b\mathbf c^\top$, cyclicity of the trace gives
$\lVert\mathbf H\rVert^2_{\mathsf G'}=(\mathbf b^\top\mathsf G'\mathbf b)(\mathbf c^\top\mathsf G'^{-1}\mathbf c)$,
and likewise for $\mathsf G''$. If $\mathsf G''=\mathrm c_{\mathrm s}\mathsf G'$ with $\mathrm c_{\mathrm s}>0$,
the factors $\mathrm c_{\mathrm s}^{-1}$ and $\mathrm c_{\mathrm s}$ in \cref{eq:link-trace-form} cancel for
every $\mathbf H$. Conversely, fixing $\mathbf c\neq0$ gives
$\mathbf b^\top\mathsf G''\mathbf b=\mathrm c_{\mathrm s}\mathbf b^\top\mathsf G'\mathbf b$ for every
$\mathbf b$, with
$\mathrm c_{\mathrm s}=(\mathbf c^\top\mathsf G'^{-1}\mathbf c)/(\mathbf c^\top\mathsf G''^{-1}\mathbf c)>0$,
and symmetric matrices with proportional quadratic forms are proportional. \hfill$\square$

\propositionheading{Properties of the loop defect}{prop:link-loop-defect}
In \Cref{def:link-loop-defect}, let $p=(e_0,\ldots,e_{r-1})$ be a closed walk based at $i_0$.

(a) The form \cref{eq:link-trace-form} is an inner product on $\mathfrak{gl}(\mathbb V)$ with
$\lVert\mathbf H\mathbf H'\rVert_{\mathsf G}\leq\lVert\mathbf H\rVert_{\mathsf G}\lVert\mathbf H'\rVert_{\mathsf G}$
and $\lVert I_{\mathbb V}\rVert_{\mathsf G}=(\dim\mathbb V)^{1/2}$. The defect $\mathfrak Q^x_p$ is a
continuous function on $\mathscr T$ that depends on $\mathring\omega^x$ only through the links; it is
nonnegative and vanishes exactly when $\rho_x(\mathring\Theta^x_p)=I_{\mathbb V}$, equivalently when
$\rho_x(\mathrm{Hol}^x(p))=I_{\mathbb V}$.

(b) Passive re-expressions leave $\mathfrak Q^x_p$ unchanged. Under the reference change
$\sigma'=\sigma\cdot g_0$, with $\mathbf R=\rho_x(g_0)$ and the generators
$\operatorname{Ad}_{g_0^{-1}}\mathring\omega^x$ of \Cref{prop:link-generator-gauge}, the defect computed
at $\sigma'$ with the matrix $\mathbf R^\top\mathsf G^{\mathrm{fib}}\mathbf R$, which represents the same
inner product in the coordinates of $\sigma'$, equals $\mathfrak Q^x_p$ for every link family, while the
defect computed at $\sigma'$ with the matrix $\mathsf G^{\mathrm{fib}}$ equals the defect at $\sigma$
computed with $\mathbf R^{-\top}\mathsf G^{\mathrm{fib}}\mathbf R^{-1}$. The latter equals
$\mathfrak Q^x_p$ for every generator family whenever $\mathbf R^\top\mathsf G^{\mathrm{fib}}\mathbf R$ is
a positive multiple of $\mathsf G^{\mathrm{fib}}$, that is, whenever $\mathbf R$ is a nonzero multiple of a
$\mathsf G^{\mathrm{fib}}$-isometry, and, when
$\mathrm d\rho_x(\mathfrak g)=\mathfrak{gl}(\mathbb V)$ and $p$ traverses some reversal pair exactly
once, only then. Likewise, the defect of the rotated walk $(e_1,\ldots,e_{r-1},e_0)$ based at $i_1$ is
that of $p$ computed with
$\rho_x(\mathring\Theta^x_{e_0})^{-\top}\mathsf G^{\mathrm{fib}}\rho_x(\mathring\Theta^x_{e_0})^{-1}$.

(c) For each agent $i$ let
$\mathsf G^{\mathrm{fib}}_i=\rho_x(U_i)^{-\top}\mathsf G^{\mathrm{fib}}\rho_x(U_i)^{-1}$, the matrix of
the same inner product in the coordinates of $u_i$. Under passive re-expression
$\mathsf G^{\mathrm{fib}}_i$ becomes $(R^x_i)^{-\top}\mathsf G^{\mathrm{fib}}_i(R^x_i)^{-1}$, the law of
\Cref{rem:geo-loop-penalties}; the matrices of different agents satisfy
$\mathsf G^{\mathrm{fib}}_i=\rho_x(T_{ij})^{-\top}\mathsf G^{\mathrm{fib}}_j\rho_x(T_{ij})^{-1}$; and, with
$\mathsf U=\rho_x(\mathrm{Hol}^x(p))$,
\begin{equation}
\mathfrak Q^x_p=\tfrac12\Tr\bigl((\mathsf G^{\mathrm{fib}}_{i_0})^{-1}(I_{\mathbb V}-\mathsf U)^\top\mathsf G^{\mathrm{fib}}_{i_0}(I_{\mathbb V}-\mathsf U)\bigr),
\label{eq:link-defect-local}
\end{equation}
which is one half of the invariant penalty of \Cref{rem:geo-loop-penalties} at $i_0$. Conversely, if
matrices $\widetilde{\mathsf G}_{i_0}$ assigned to the frames of agent $i_0$ transform by the law of that
remark, then $\rho_x(U_{i_0})^\top\widetilde{\mathsf G}_{i_0}\rho_x(U_{i_0})$ is unchanged by passive
re-expression, and $\widetilde{\mathsf G}_{i_0}$ is its frame image.

(d) If $\rho_x(\mathring\Theta^x_p)$ is a $\mathsf G^{\mathrm{fib}}$-isometry, then
$\mathfrak Q^x_p=\dim\mathbb V-\Tr\rho_x(\mathring\Theta^x_p)=\dim\mathbb V-\Tr\mathsf U\in[0,2\dim\mathbb V]$,
which is $\dim\mathbb V$ times the plaquette term $1-\Tr\mathsf U/\dim\mathbb V$ of Wilson's lattice
action \citep[Eq.~(3.29)]{WilsonConfinement1974}, in the normalization of
\citet[Ch.~7, Eq.~(7.6)]{Creutz1983}.

(e) Let $\mathring\Omega_a=\mathrm d\rho_x(\mathring\omega^x_{e_a})$ and
$\mathring\varsigma_p=\sum_{a=0}^{r-1}\lVert\mathring\Omega_a\rVert_{\mathsf G}$. Then
\begin{equation}
\Bigl\lvert\mathfrak Q^x_p-\tfrac12\bigl\lVert\mathrm d\rho_x(\mathbf C\mathring\omega^x)_p\bigr\rVert^2_{\mathsf G}\Bigr\rvert
\leq\mathring\varsigma_p\bigl(\exp(\mathring\varsigma_p)-1-\mathring\varsigma_p\bigr)
+\tfrac12\bigl(\exp(\mathring\varsigma_p)-1-\mathring\varsigma_p\bigr)^2,
\label{eq:link-defect-expansion}
\end{equation}
and the right side is at most
$(\tfrac12\exp(1)+\tfrac18\exp(2))\mathring\varsigma_p^3<\tfrac{23}{10}\mathring\varsigma_p^3$ when
$\mathring\varsigma_p\leq1$. When $\mathrm d\rho_x$ is injective and $\mathscr T$ carries the norm
$\lVert\mathring\omega\rVert^2=\sum_{e\in E_+}\lVert\mathring\omega_e\rVert^2_{\mathsf G}$, one has
$\mathring\varsigma_p\leq r\lVert\mathring\omega^x\rVert$, so the bound is
$\tfrac{23}{10}r^3\lVert\mathring\omega^x\rVert^3$ on the ball $r\lVert\mathring\omega^x\rVert\leq1$.
\status{ESTABLISHED}

The proof is in \cref{sec:proofs-links}.

The matrix at the reference point is thus the one datum beyond the group that the defect uses. Carried
along with the reference point, it leaves the defect invariant under every reference change; held as a
matrix, it leaves the defect invariant under the reference changes acting by nonzero multiples of its
isometries, and, when $\mathrm d\rho_x(\mathfrak g)=\mathfrak{gl}(\mathbb V)$ and the walk traverses some
reversal pair exactly once, under no other (\Cref{prop:link-loop-defect}(b)). The symmetry of the measurement is reduced in this way from all
reference changes to those acting by multiples of isometries of $\mathsf G^{\mathrm{fib}}$; on
holonomies that are isometries the defect is Wilson's trace density (part (d)), and one reference matrix
supplies the metric of \Cref{rem:geo-loop-penalties} at every vertex and for every closed walk (part
(c)). A cyclic rotation of a closed walk replaces the matrix by its image under the twist moved across
(part (b)), so the base vertex does not affect the defect when that twist acts by a multiple of an
isometry. The quadratic part of the defect is carried by the cycle map (part (e)), which reappears in
\Cref{prop:link-linearized-equation}.

\section{The generator prior}
\label{sec:link-prior}

The generator chart supplies a link prior for every Lie group: a Gaussian base law on the generators,
which exists whether or not $G$ is compact, multiplied by the exponential of minus a nonnegative loop
action (\Cref{prop:link-generator-prior}). It replaces the Haar reference, whose normalizability the scope
of \Cref{ch:geometry} declines for noncompact and general linear $G$, by a reference on the Lie algebra,
and the abelian links of \Cref{ch:finite-depth-formation} are its translation-group instance
(\Cref{prop:link-translation-instance}).

\definitionheading{Generator prior}{def:link-generator-prior}
In \Cref{def:link-exponential-chart,def:link-loop-defect}, equip $\mathfrak g$ with an inner product,
written $\langle\cdot,\cdot\rangle_{\mathsf G}$ (the pulled-back inner product of
\Cref{def:link-loop-defect} when $\mathrm d\rho_x$ is injective), and $\mathscr T$ with the product inner
product and the Lebesgue measure $\mathrm d\mathring\omega$ normalized by it, under which the cube
spanned by an orthonormal basis has volume one. For each $e\in E_+$ let the anchor
$\mathring\alpha^x_e\in\mathfrak g$ and the covariance $\mathring\Sigma^x_e$, a positive definite
self-adjoint operator on $(\mathfrak g,\langle\cdot,\cdot\rangle_{\mathsf G})$, be generative parameters,
that is, coordinates of $\theta$, and let $\mathbb L$ be a finite family of closed walks with coefficients
$\mathrm c_p\geq0$. The base law, link action, normalizer, and generator prior are
\begin{gather}
\nu^{\mathscr T}=\bigotimes_{e\in E_+}\mathcal N(\mathring\alpha^x_e,\mathring\Sigma^x_e),
\qquad
\mathscr S(\mathring\omega^x)=\sum_{p\in\mathbb L}\mathrm c_p\mathfrak Q^x_p(\mathring\omega^x),
\label{eq:link-base-action}\\
Z_{\mathscr S}=\int_{\mathscr T}e^{-\mathscr S}d\nu^{\mathscr T},
\qquad
Z_{\mathscr S}^{-1}e^{-\mathscr S}\nu^{\mathscr T},
\label{eq:link-generator-prior}
\end{gather}
where $\mathcal N(\mathring\alpha,\mathring\Sigma)$ has the density
$(2\pi)^{-\dim\mathfrak g/2}(\det\mathring\Sigma)^{-1/2}\exp\bigl(-\tfrac12\langle\mathring\omega-\mathring\alpha,\mathring\Sigma^{-1}(\mathring\omega-\mathring\alpha)\rangle_{\mathsf G}\bigr)$
with respect to the normalized Lebesgue measure of $\mathfrak g$. With independent links in the two
channels the link space is $\mathscr T\times\mathscr T$, the base law is the product of the two base
laws, and the action is the sum of the two channel actions; under the common-link specialization one
base law on $\mathscr T$ serves both channels and the action is the sum over channels.

\propositionheading{Normalization and covariance of the generator prior}{prop:link-generator-prior}
Let $G$ be a Lie group, compact or not, and let the data of \Cref{def:link-generator-prior} be given.

(a) The action $\mathscr S$ is finite, continuous, and nonnegative on $\mathscr T$ and has finite
$\nu^{\mathscr T}$-mean, and
\begin{equation}
0<\exp\bigl(-\E_{\nu^{\mathscr T}}\mathscr S\bigr)\leq Z_{\mathscr S}\leq1;
\label{eq:link-normalizer-bounds}
\end{equation}
when $\rho_x(G)$ consists of $\mathsf G^{\mathrm{fib}}$-isometries, also
$Z_{\mathscr S}\geq\exp(-2\dim\mathbb V\sum_p\mathrm c_p)$. The generator prior is a probability law on
the standard Borel space $\mathscr T$ with a strictly positive continuous density with respect to
$\mathrm d\mathring\omega$. Hence $\nu^{\mathscr T}$, $\mathscr S$, and $Z_{\mathscr S}$ satisfy the
hypotheses that \Cref{prop:elbo-random-link-vfe} places on its base law, action, and normalizer, with the
generator family as the latent link variable and the links $\exp_\sigma(\mathring\omega^x)$ as a
continuous function of it. With $\widetilde{\mathscr S}$ the action written as a function of the link
family through $\mathring\Theta^x_p=U_{i_0}^{-1}\mathrm{Hol}^x(p)U_{i_0}$, so that
$\mathscr S=\widetilde{\mathscr S}\circ\exp_\sigma$, the image of the prior under $\exp_\sigma$ is
$Z_{\mathscr S}^{-1}e^{-\widetilde{\mathscr S}}(\exp_\sigma)_\#\nu^{\mathscr T}$, with the same normalizer.

(b) Passive re-expressions leave the generator prior unchanged and carry its image on link families by
the vertex law \cref{eq:geo-regime-two-gauge-law}.

(c) Let $\sigma'=\sigma\cdot g_0$ and $\mathbf R=\rho_x(g_0)$, and let $\operatorname{Ad}_{g_0^{-1}}$ act
on every component of $\mathscr T$. The image of the generator prior under $\operatorname{Ad}_{g_0^{-1}}$
is the generator prior relative to $\sigma'$ with the anchors
$\operatorname{Ad}_{g_0^{-1}}\mathring\alpha^x_e$, the covariances
$\operatorname{Ad}_{g_0^{-1}}\mathring\Sigma^x_e\operatorname{Ad}_{g_0}$, the fiber matrix
$\mathbf R^\top\mathsf G^{\mathrm{fib}}\mathbf R$, the inner product
$\langle\operatorname{Ad}_{g_0}\cdot,\operatorname{Ad}_{g_0}\cdot\rangle_{\mathsf G}$ on $\mathfrak g$, which is the
pullback for $\mathbf R^\top\mathsf G^{\mathrm{fib}}\mathbf R$ when $\mathrm d\rho_x$ is injective, and the
same walks and coefficients; both priors have the same image on link families.

(d) If $\operatorname{Ad}_{g_0}$ is an isometry of $\langle\cdot,\cdot\rangle_{\mathsf G}$, fixes every
anchor, and commutes with every covariance, and $\mathbf R$ is a nonzero multiple of a
$\mathsf G^{\mathrm{fib}}$-isometry, then the image of the generator prior under
$\operatorname{Ad}_{g_0^{-1}}$ is the generator prior itself, so the reference change by $g_0$ leaves the
anchors, covariances, and coefficients unchanged and multiplies the fiber matrix by a positive scalar,
which changes neither the inner product nor the defect. For the pulled-back inner product the isometry
condition follows from the condition on $\mathbf R$, and when $\mathrm d\rho_x(\mathfrak g)=\mathfrak{gl}(\mathbb V)$
it is equivalent to it. \status{ESTABLISHED}

The proof is in \cref{sec:proofs-links}.

The Haar-referenced density of \Cref{rem:geo-loop-penalties} needs the finite invariant measure of a
compact group, whereas the base law of \Cref{def:link-generator-prior} lives on the Lie algebra, so
\cref{eq:link-normalizer-bounds} holds for every Lie group; the price is a set of anchors and covariances
that a reference change transforms (\Cref{prop:link-generator-prior}).

\propositionheading{Translation links}{prop:link-translation-instance}
Let $G=(\R^K,+)$ act on the sample fiber $\R^K$ by translation, and write the group additively.

(a) The affine representation and its differential,
\begin{equation}
\rho(\mathring\omega)=\begin{pmatrix}I_K&\mathring\omega\\0&1\end{pmatrix},
\qquad
\mathrm d\rho(\mathring\omega)=\begin{pmatrix}0&\mathring\omega\\0&0\end{pmatrix},
\label{eq:link-affine-representation}
\end{equation}
form a faithful smooth linear representation on $\R^{K+1}$ with $\mathfrak g=\R^K$ and $\exp$ the
identity, and $\rho(\mathring\omega)(k,1)=(k+\mathring\omega,1)$, so the translations of the sample fiber
act on the affine slice $\R^K\times\{1\}$. A link $\Theta_e\in\R^K$ on $e:j\to i$, acting by
$k\mapsto k+\Theta_e$, has the generator $\mathring\omega_e=\Theta_e-U_i+U_j$, and the reference holonomy
of a closed walk is the loop sum $\mathring\Theta_p=(\mathbf C\mathring\omega)_p=\sum_a\Theta_{e_a}$, in
which the frames cancel.

(b) For every symmetric positive definite $\mathsf G^{\mathrm{fib}}$ on $\R^{K+1}$, the pulled-back norm
and the loop defect are
\begin{equation}
\lVert\mathring\omega\rVert^2_{\mathsf G}
=\bigl((\mathsf G^{\mathrm{fib}})^{-1}\bigr)_{K+1,K+1}(\mathring\omega,0)^\top\mathsf G^{\mathrm{fib}}(\mathring\omega,0),
\qquad
\mathfrak Q_p=\tfrac12\bigl\lVert(\mathbf C\mathring\omega)_p\bigr\rVert^2_{\mathsf G},
\label{eq:link-translation-defect}
\end{equation}
so the left side of \cref{eq:link-defect-expansion} vanishes; for $\mathsf G^{\mathrm{fib}}=I_{K+1}$,
$\mathfrak Q_p=\tfrac12\lvert\sum_a\Theta_{e_a}\rvert^2$.

(c) In the four-cycle experiment of \Cref{def:rg-formation-cycle-experiment} and the gated translation
family of \Cref{def:rg-gated-translation-family}, where $K=1$ and the link translations $\omega_e$ are
the links $\Theta_e$, the configuration prior is $\mathcal N(0,I)$, so the link translations are
independent of the means and have the law $\prod_e\mathcal N(0,1)$. In generator coordinates this is the
base law of \Cref{def:link-generator-prior} for $\mathsf G^{\mathrm{fib}}=I_2$ with unit covariances and
the anchors $\mathring\alpha_e=U_j-U_i$, which vanish when every frame equals $\sigma$ and always have
zero loop sums, and the action is $\mathscr S=0$. The loop $p=(e_3,e_2,e_1,e_0)$ of the four-cycle,
based at agent $0$, has $(\mathbf C\mathring\omega)_p=\sum_i\omega_i$ and
$\mathfrak Q_p=\tfrac12(\sum_i\omega_i)^2$; in these models it enters only through the edge costs, whose
residuals satisfy $\sum_ir^b_i=\sum_ir^m_i=-\sum_i\omega_i$.

(d) In the triangle-chain experiment of \Cref{def:rg-formation-experiment}, with the symbols of
\cref{tab:rg-formation-notation} and $\mathsf G^{\mathrm{fib}}=I_2$, every cyclically oriented triangle
$q$ has $\mathfrak h_q=(\mathbf C\mathring\omega)_q$ and loop energy
$\mathrm J_1h_q^2=\tfrac23\mathrm J_1\mathfrak Q_q$, and for every configuration
\begin{equation}
\frac{\mathbf x^\top\Omega_0(\vartheta)\mathbf x}{2\mathrm T}
=\sum_q\frac{\vartheta\mathrm J_1}{3\mathrm T}\mathfrak Q_q+\frac{\mathbf x^\top\Omega^\perp_0(\vartheta)\mathbf x}{2\mathrm T},
\qquad
\Omega^\perp_0(\vartheta)=\vartheta B^\top\operatorname{diag}(\mathrm J_\ell I)B+P_{\mathcal K}-\mathrm TR_{\mathrm{ev}},
\label{eq:link-triangle-chain-prior}
\end{equation}
where the kernel of $\Omega^\perp_0(\vartheta)$ is the loop space $\mathcal H$ and $\Omega^\perp_0(\vartheta)$ is
positive definite on $\mathcal U\oplus\mathcal K$. Under the configuration prior the loop sums
$\mathfrak h_q$ are therefore independent with law $\mathcal N(0,3\mathrm T/(\vartheta\mathrm J_1))$ and
independent of the projection of $\mathbf x$ onto $\mathcal U\oplus\mathcal K$, and the quadratic energies
$\mathrm E_\vartheta$ and $\mathrm E^{\mathrm{phys}}_\vartheta$ of
\cref{eq:rg-formation-partition-integrals} are $\tfrac13\vartheta\mathrm J_1\sum_q\mathfrak Q_q$ plus a
function of that projection. This prior differs from \Cref{def:link-generator-prior} in that its second
factor is a Gaussian energy of means and links jointly that is constant along $\mathcal H$, so it is not
integrable alone, and the loop defect with coefficient $\vartheta\mathrm J_1/(3\mathrm T)$ supplies all
of the prior's precision along loop coordinates. \status{ESTABLISHED}

The proof is in \cref{sec:proofs-links}.

\remarkheading{The running example in the chart}{rem:link-running-example}
In the running example of \Cref{sec:geo-running-example} take $E_+=\{e_1,e_2,e_3\}$, $\sigma=u_1$, and
$U_i=e_G$, and write $R_x(\phi)=\exp(\phi(E_{32}-E_{23}))$ and $R_z(\phi)=\exp(\phi(E_{21}-E_{12}))$ for
the rotations about the first and third axes. The chart covers every link family of $SO(3)$
(\Cref{prop:link-chart-image}), and the generators of rotation angle at most $\pi$ are
$\mathring\omega^b_{e_1}=\tfrac\pi2(E_{21}-E_{12})$, $\mathring\omega^b_{e_2}=\tfrac\pi2(E_{32}-E_{23})$,
and $\mathring\omega^b_{e_3}=0$. The reference holonomy of $p=(e_1,e_2,e_3)$ is
$\mathring\Theta^b_p=R_z(\tfrac\pi2)R_x(\tfrac\pi2)$, the cyclic permutation matrix that sends each
standard basis vector to the next one, and it equals $\mathrm{Hol}^b(p)$, a rotation by $2\pi/3$ about the
axis spanned by $(1,1,1)$. With $\mathsf G^{\mathrm{fib}}=I_3$ every rotation is an isometry, a rotation
by $\phi$ has defect $4\sin^2(\phi/2)\in[0,4]$, and $\mathfrak Q^b_p=3-\Tr\mathring\Theta^b_p=3$, while the
trivial model representation gives $\mathfrak Q^m_p=0$. The cycle map gives
$(\mathbf C\mathring\omega^b)_p=\tfrac\pi2(E_{21}-E_{12}+E_{32}-E_{23})$, whose quadratic term
$\tfrac12\lVert(\mathbf C\mathring\omega^b)_p\rVert^2_{\mathsf G}=\pi^2/2$ exceeds the defect; here
$\mathring\varsigma_p=\pi\sqrt2>1$, outside the ball of \Cref{prop:link-loop-defect}(e). A passive
re-expression $g_2$ changes $\Theta^b_{e_1}$ and $\Theta^b_{e_2}$ but no generator.
\status{ESTABLISHED}

\noindent\emph{Proof.}
The series gives $R_z(\tfrac\pi2)$ and $R_x(\tfrac\pi2)$ as the quarter turns of \cref{eq:geo-running-links},
whose product is the displayed permutation matrix; it fixes $(1,1,1)$ and has trace $0=1+2\cos\phi$ with
$\phi=2\pi/3$. A rotation by $\phi$ has eigenvalues $1$ and $e^{\pm i\phi}$, so
\Cref{prop:link-loop-defect}(d) gives the defect $3-(1+2\cos\phi)=4\sin^2(\phi/2)$, and
$\mathfrak Q^b_p=3-0$. The four nonzero entries of $(\mathbf C\mathring\omega^b)_p$ have modulus $\pi/2$,
so its squared Frobenius norm is $\pi^2$, and each nonzero generator has norm $\pi/\sqrt2$. Generators
are unchanged by passive re-expression by \Cref{prop:link-generator-gauge}(a). \hfill$\square$

\section{Latent generators in the collective evidence bound}
\label{sec:link-latent}

What joint law and what collective free energy result when the link values of a channel leave the
structural data and the generators of the exponential chart become latent coordinates? The joint has
unit mass on every finite record hypergraph, and its free energy splits into a link sector and the
conditional free energies of the held-link model (\cref{prop:link-latent-normalization}); for a Gaussian
recognition law of the generators every term is finite and the link sector is the relative entropy from
the generator prior (\cref{prop:link-gaussian-recognition-vfe}). Throughout, the generator prior is that
of \cref{def:link-generator-prior} in the following form: the base law $\nu^{\mathscr T}$ is the product,
over the copies $e\in E_+$, of the Gaussian laws
$\mathcal N(\mathring\alpha^x_e,\mathring\Sigma^x_e)$ on $(\mathfrak g,\langle\cdot,\cdot\rangle_{\mathsf G})$,
written $\mathcal N(\mathring\alpha,\mathring\Sigma)$ with a block-diagonal covariance $\mathring\Sigma$; the
action $\mathscr S$ is a finite nonnegative combination of the loop defects $\mathfrak Q^x_p$ of
\cref{def:link-loop-defect}; and $0<Z_{\mathscr S}\leq1$ by \cref{prop:link-generator-prior}. The
definitions and statements of this section and of \cref{sec:link-field,sec:link-held,sec:link-induced} are
written for the link space $\mathscr T$ of one latent channel $x$, which covers the common-link
specialization. When both channels carry independent latent generator families, the link space is
$\mathscr T\times\mathscr T$ with the product of the two base laws and the sum of the two actions
(\cref{def:link-exponential-chart,def:link-generator-prior}); every definition and statement then holds with
these data in place of $\mathscr T$, $\nu^{\mathscr T}$, and $\mathscr S$, with the summands indexed by the
pairs $(x,e)$, by the same proof, since the proofs use only that the link space is a finite-dimensional
inner product space carrying a Gaussian base law that factors over orthogonal summands and an action that
is a finite nonnegative combination of loop defects. The space
$\mathscr T$ carries the orthogonal sum of the inner products of its summands, its norm
$\lVert\cdot\rVert$, and the Lebesgue measure normalized by it; gradients and Hessians on $\mathscr T$
refer to this inner product, $\lVert\cdot\rVert$ is the operator norm on matrices, and upright
$\mathrm c_1,\mathrm c_2,\ldots$ are constants local to one statement.

\definitionheading{Joint law with latent generators}{def:link-latent-model}
Let the structural data $X$ contain the interaction multigraph $\mathfrak I$, an orientation set $E_+$ with
one copy of each reversal pair, the reference point $\sigma$ with the frames $U_i$, $u_i=\sigma\cdot
U_i^{-1}$, and the acyclic graph $\Gamma$ of \cref{def:gen-parameter-split}, but no link values of the
channels $x$ whose links are latent. For those channels the generator family
$\mathring\omega=(\mathring\omega^x_e)$ is a latent coordinate with values in $\mathscr T$, or in
$\mathscr T\times\mathscr T$ when both channels are latent, and it determines the links of the exponential
chart of
\cref{def:link-exponential-chart},
\begin{equation}
\Theta^x_e(\mathring\omega)=U_i\exp(\mathring\omega^x_e)U_j^{-1},
\qquad
\Theta^x_{\bar e}(\mathring\omega)=\Theta^x_e(\mathring\omega)^{-1},
\qquad e:j\to i\text{ in }E_+,
\label{eq:link-generator-links}
\end{equation}
a continuous map $\mathring\omega\mapsto\Theta(\mathring\omega)$ into the link families. Write $(X,\Theta)$ for
$X$ with a link family $\Theta$ adjoined, and let the factors \eqref{eq:gen-kernels} and the record kernels
of \cref{hyp:local-interaction-kernels} be measurable in the link values, so that for each $\Theta$ the
held-link model of \cref{ch:generative,ch:elbo} at $(X,\Theta)$ has the joint
$\mathbb P_{\theta,V}(dy,do\given X,\Theta)=\mathbb P^{\mathsf Y}_{\theta,V}(dy\given X,\Theta)\prod_{a\in\mathfrak A}K_a(X,\Theta,y_{\partial a},do_a)$,
the evidence $p_{\theta,V}(o\given X,\Theta)$, the posterior $\boldsymbol\Pi_{\theta,V,o,X,\Theta}$, and, for
laws $Q$ on $\mathsf Y$, the free energy
$\Fenergy_V[Q;X,\Theta,o]=-\log p_{\theta,V}(o\given X,\Theta)+\KL(Q\Vert\boldsymbol\Pi_{\theta,V,o,X,\Theta})$ of
\cref{def:elbo-extended}. With the generator prior
$\mathbb P^{\mathscr T}_\theta(d\mathring\omega\given X)=Z_{\mathscr S}^{-1}e^{-\mathscr S(\mathring\omega)}\nu^{\mathscr T}(d\mathring\omega)$,
the joint law with latent generators is the law on $\mathscr T\times\mathsf Y\times\mathsf O$
\begin{equation}
\mathbb P_{\theta,V}(d\mathring\omega,dy,do\given X)
=\mathbb P^{\mathscr T}_\theta(d\mathring\omega\given X)
\mathbb P^{\mathsf Y}_{\theta,V}\bigl(dy\given X,\Theta(\mathring\omega)\bigr)
\prod_{a\in\mathfrak A}K_a\bigl(X,\Theta(\mathring\omega),y_{\partial a},do_a\bigr),
\label{eq:link-latent-joint}
\end{equation}
whose evidence $p_{\theta,V}(o\given X)$, posterior $\boldsymbol\Pi_{\theta,V,o,X}(d\mathring\omega,dy)$, and
extended free energy $\Fenergy_V[Q;X,o]$ for laws $Q$ on $\mathscr T\times\mathsf Y$ are those of
\cref{def:elbo-extended} with the latent space $\mathscr T\times\mathsf Y$ and the baseline
$\mathbb P^{\mathscr T}_\theta(d\mathring\omega\given X)\mathbb P^{\mathsf Y}_{\theta,V}(dy\given X,\Theta(\mathring\omega))$.
With $\theta$, $X$ and $o$ held, write
$p(o\given\mathring\omega)=p_{\theta,V}(o\given X,\Theta(\mathring\omega))$,
$\boldsymbol\Pi_{\mathring\omega}=\boldsymbol\Pi_{\theta,V,o,X,\Theta(\mathring\omega)}$,
$\mathbb P^{\mathsf Y}(\cdot\given\mathring\omega)=\mathbb P^{\mathsf Y}_{\theta,V}(\cdot\given X,\Theta(\mathring\omega))$,
and, for records with densities \eqref{eq:obs-interaction-density},
$L_o(\mathring\omega,y)=\prod_{a}e^{-E_{a,o}(X,\Theta(\mathring\omega),y_{\partial a})}$ and
$H_o(\mathring\omega,y)=\sum_aE_{a,o}(X,\Theta(\mathring\omega),y_{\partial a})$. The anchors
$\mathring\alpha^x_e$, the prior covariances $\mathring\Sigma^x_e$, and the loop coefficients are coordinates
of $\theta$, and the walks of $\mathscr S$ belong to $X$. The prior reads $\theta$ and $X$; each agent factor
reads $\theta$, $X$, its parents and $\mathring\omega$; each record reads $\theta$, $X$, $y_{\partial a}$ and
$\mathring\omega$. No factor reads a recognition law, a variational parameter, or the posterior, so
\cref{req:gen-typing-prohibition} holds: moving an anchor or a prior covariance is an M-coordinate
(\cref{def:elbo-generalized-m-acceptance}), while the generators are moved by recognition. Passive
re-expression changes the frames and no generator (\cref{prop:geo-edge-twist-invariance}); a change of
reference point acts on the generator prior as in \cref{prop:link-generator-gauge}.

\propositionheading{Normalization, evidence, and the conditional decomposition}{prop:link-latent-normalization}
Assume \cref{def:link-latent-model}. (a) For every finite record hypergraph, cyclic or not, and every
$(\theta,X)$, \eqref{eq:link-latent-joint} is a probability law on $\mathscr T\times\mathsf Y\times\mathsf O$,
and the generator family is a root block of the directed latent law. (b) If the records have densities
\eqref{eq:obs-interaction-density}, then
$\mathring\omega\mapsto p(o\given\mathring\omega)=\int L_o(\mathring\omega,y)\mathbb P^{\mathsf Y}(dy\given\mathring\omega)$
is measurable and
\begin{equation}
p_{\theta,V}(o\given X)=\int_{\mathscr T}p(o\given\mathring\omega)\mathbb P^{\mathscr T}_\theta(d\mathring\omega\given X);
\label{eq:link-latent-evidence}
\end{equation}
the evidence is positive exactly when $p(o\given\cdot)$ is positive on a set of positive Lebesgue measure,
and if every record density is bounded, $e^{-E_{a,o}}\leq\mathrm c_a$, then
$p(o\given\mathring\omega)\leq\mathrm c_{\mathfrak A}:=\prod_a\mathrm c_a$ for every $\mathring\omega$ and the evidence is at
most $\mathrm c_{\mathfrak A}$; binary success records have $\mathrm c_a=1$. (c) Suppose further that the record
densities are bounded and $p(o\given\mathring\omega)>0$ for every $\mathring\omega$. Then the posterior
disintegrates as
$\boldsymbol\Pi_{\theta,V,o,X}(d\mathring\omega,dy)=\boldsymbol\Pi^{\mathscr T}(d\mathring\omega)\boldsymbol\Pi_{\mathring\omega}(dy)$
with link marginal
$\boldsymbol\Pi^{\mathscr T}(d\mathring\omega)=p_{\theta,V}(o\given X)^{-1}p(o\given\mathring\omega)\mathbb P^{\mathscr T}_\theta(d\mathring\omega\given X)$,
and every law $Q$ on $\mathscr T\times\mathsf Y$, disintegrated as $Q(d\mathring\omega,dy)=R(d\mathring\omega)Q_{\mathring\omega}(dy)$,
satisfies
\begin{equation}
\Fenergy_V[Q;X,o]
=\KL\bigl(R\Vert\mathbb P^{\mathscr T}_\theta(\cdot\given X)\bigr)
+\int_{\mathscr T}\Fenergy_V\bigl[Q_{\mathring\omega};X,\Theta(\mathring\omega),o\bigr]R(d\mathring\omega)
\qquad\text{in }(-\infty,+\infty],
\label{eq:link-decomposition}
\end{equation}
with a measurable integrand bounded below by $-\log\mathrm c_{\mathfrak A}$. When $\KL(R\Vert\nu^{\mathscr T})<\infty$
and $\E_R\lvert\mathscr S\rvert<\infty$, the first term equals
$\KL(R\Vert\nu^{\mathscr T})+\E_R\mathscr S+\log Z_{\mathscr S}$, and \eqref{eq:link-decomposition} is
\eqref{eq:rg-geometry-random-link-vfe} for the latent variable $\mathring\omega$. \status{ESTABLISHED}

The proof is in \cref{sec:proofs-links}.

\lemmaheading{Exponential moments and growth of loop actions}{lem:link-growth}
(a) For a Gaussian law $\mathcal N(\mu,\Sigma)$ on $\mathscr T$ and every $\mathrm c\geq0$,
$\E_{\mathcal N(\mu,\Sigma)}e^{\mathrm c\lVert\mathring\omega\rVert}\leq e^{\mathrm c\lVert\mu\rVert}\bigl(2e^{\mathrm c^2\lVert\Sigma\rVert/2}\bigr)^{\dim\mathscr T}$,
and the same bound holds on each summand of $\mathscr T$ with its dimension.
(b) For a representation $\rho_x$ on a real vector space of finite dimension, the map
$\mathring\xi\mapsto\rho_x(\exp\mathring\xi)=\exp(\mathrm d\rho_x\mathring\xi)$ on
$(\mathfrak g,\langle\cdot,\cdot\rangle_{\mathsf G})$ is smooth, and its value and its first and second
derivatives at $\mathring\xi$ have norms at most $e^{\lVert\mathrm d\rho_x\rVert\lVert\mathring\xi\rVert}$,
$\lVert\mathrm d\rho_x\rVert e^{\lVert\mathrm d\rho_x\rVert\lVert\mathring\xi\rVert}$ and
$\lVert\mathrm d\rho_x\rVert^2e^{\lVert\mathrm d\rho_x\rVert\lVert\mathring\xi\rVert}$. (c) Every real polynomial
$\mathbf h$ in the matrix entries of the represented twists
$\rho_x(\mathring\Theta^x_e)^{\pm1}=\exp(\pm\mathrm d\rho_x\mathring\omega^x_e)$, $e\in E_+$, is smooth on
$\mathscr T$ and satisfies, for constants depending on $\mathbf h$,
\begin{equation}
\lvert\mathbf h(\mathring\omega)\rvert+\lVert\nabla\mathbf h(\mathring\omega)\rVert
+\lVert\nabla^2\mathbf h(\mathring\omega)\rVert\leq\mathrm c_1e^{\mathrm c_2\lVert\mathring\omega\rVert};
\label{eq:link-growth}
\end{equation}
each loop defect of \cref{def:link-loop-defect}, and hence the action $\mathscr S$, is such a polynomial.
\status{ESTABLISHED}

\noindent\emph{Proof.}
(a) Write $\mathring\omega=\mu+\Sigma^{1/2}\mathbf b$ with $\mathbf b$ standard Gaussian in an orthonormal
basis, so that $\lVert\mathring\omega\rVert\leq\lVert\mu\rVert+\lVert\Sigma\rVert^{1/2}\sum_l\lvert\mathbf b_l\rvert$;
the coordinates are independent and, for $\mathrm c'\geq0$,
$\E e^{\mathrm c'\lvert\mathbf b_l\rvert}\leq\E e^{\mathrm c'\mathbf b_l}+\E e^{-\mathrm c'\mathbf b_l}=2e^{\mathrm c'^2/2}$.
(b) The identity $\rho_x(\exp\mathring\xi)=\exp(\mathrm d\rho_x\mathring\xi)$ is
\cref{lem:link-exp-naturality}. The series $\exp\mathbf B=\sum_l\mathbf B^l/l!$ converges
with all its derivatives. The first derivative of $\mathbf B\mapsto\mathbf B^l$ in the direction $\mathbf T$ is a
sum of $l$ products of $l-1$ factors $\mathbf B$ and one factor $\mathbf T$, and its second derivative in the
directions $\mathbf T$ and $\mathbf T'$ is a sum of $l(l-1)$ products of $l-2$ factors $\mathbf B$ and one factor
each of $\mathbf T$ and $\mathbf T'$. Summing the series, the value and the first and second derivatives of the
exponential have norms at most $e^{\lVert\mathbf B\rVert}$, $e^{\lVert\mathbf B\rVert}\lVert\mathbf T\rVert$ and
$e^{\lVert\mathbf B\rVert}\lVert\mathbf T\rVert\lVert\mathbf T'\rVert$; composing with the linear map $\mathrm d\rho_x$
gives (b). (c) By (b) every entry of
$\exp(\pm\mathrm d\rho_x\mathring\omega^x_e)$ and its first two derivatives are bounded by
$\mathrm c e^{\lVert\mathrm d\rho_x\rVert\lVert\mathring\omega\rVert}$, and the Leibniz rule bounds a monomial of
degree $\mathrm c_0$ and its derivatives by $\mathrm c'e^{\mathrm c_0\lVert\mathrm d\rho_x\rVert\lVert\mathring\omega\rVert}$. The
reference-coordinate holonomy of a closed walk is the product of the matrices
$\rho_x(\mathring\Theta^x_e)^{\pm1}$ along the walk (\cref{prop:geo-edge-twist-invariance}), and a loop
defect is a quadratic polynomial in its entries (\cref{def:link-loop-defect}). \hfill$\square$

\propositionheading{Gaussian recognition of the generators}{prop:link-gaussian-recognition-vfe}
Let $R=\mathcal N(\mathring\mu,\mathring\Lambda^{-1})$ be a Gaussian law on $\mathscr T$ with mean
$\mathring\mu$ and positive-definite self-adjoint precision $\mathring\Lambda$; in the mean-field case
$R=\bigotimes\mathcal N(\mathring\mu^x_e,(\mathring\Lambda^x_e)^{-1})$ over the summands of $\mathscr T$.
(a) The relative entropy from the base law is finite and equals
\begin{equation}
\KL(R\Vert\nu^{\mathscr T})
=\tfrac12\Bigl[\Tr\bigl(\mathring\Sigma^{-1}\mathring\Lambda^{-1}\bigr)
+\bigl\langle\mathring\mu-\mathring\alpha,\mathring\Sigma^{-1}(\mathring\mu-\mathring\alpha)\bigr\rangle
-\dim\mathscr T+\log\det\mathring\Sigma+\log\det\mathring\Lambda\Bigr];
\label{eq:link-gaussian-kl}
\end{equation}
in the mean-field case it is the sum of the same expressions over the summands, and in general it equals
$\operatorname{TC}(R)+\sum\KL(R^x_e\Vert\mathcal N(\mathring\alpha^x_e,\mathring\Sigma^x_e))$, with $R^x_e$ the
summand marginals and
$\operatorname{TC}(R)=\frac12[\sum\log\det(\mathring\Lambda^{-1})^x_e+\log\det\mathring\Lambda]$, where
$(\mathring\Lambda^{-1})^x_e$ are the diagonal blocks of $\mathring\Lambda^{-1}$. (b)
$\E_Re^{\mathrm c\lVert\mathring\omega\rVert}<\infty$ for every $\mathrm c\geq0$ and $\E_R\mathscr S<\infty$.
Consequently, if the record densities are bounded and $p(o\given\mathring\omega)>0$ for every
$\mathring\omega$, then for every Markov kernel $\mathring\omega\mapsto Q_{\mathring\omega}$ the free energy of
$Q=R(d\mathring\omega)Q_{\mathring\omega}(dy)$ is \eqref{eq:rg-geometry-random-link-vfe} for the latent variable
$\mathring\omega$, with the finite link sector
\begin{equation}
\KL(R\Vert\nu^{\mathscr T})+\E_R\mathscr S+\log Z_{\mathscr S}=\KL\bigl(R\Vert\mathbb P^{\mathscr T}_\theta(\cdot\given X)\bigr).
\label{eq:link-sector}
\end{equation}
(c) If moreover $Q_{\mathring\omega}\ll\mathbb P^{\mathsf Y}(\cdot\given\mathring\omega)$ for $R$-almost every
$\mathring\omega$, $\int\KL(Q_{\mathring\omega}\Vert\mathbb P^{\mathsf Y}(\cdot\given\mathring\omega))R(d\mathring\omega)<\infty$,
and $\E_Q\lvert H_o\rvert<\infty$, then the free energy is finite and equals
\begin{equation}
\Fenergy_V[Q;X,o]=\KL\bigl(R\Vert\mathbb P^{\mathscr T}_\theta(\cdot\given X)\bigr)
+\int\KL\bigl(Q_{\mathring\omega}\Vert\mathbb P^{\mathsf Y}(\cdot\given\mathring\omega)\bigr)R(d\mathring\omega)
+\E_QH_o.
\label{eq:link-gaussian-energy-form}
\end{equation}
The condition on $H_o$ holds for a product law $Q=R\otimes Q^{\mathsf Y}$ when $\mathsf Y$ is a real vector
space of finite dimension, $Q^{\mathsf Y}$ has finite second moments, and every record satisfies
$\lvert E_{a,o}(X,\Theta(\mathring\omega),y_{\partial a})\rvert\leq\mathrm c_1(1+\lVert y\rVert^2)e^{\mathrm c_2\lVert\mathring\omega\rVert}$.
\status{ESTABLISHED}

\noindent\emph{Proof.}
(a) Both laws have densities with respect to Lebesgue measure on $\mathscr T$, and for a self-adjoint
operator $\mathbf B$, $\E_R\langle\mathring\omega-\mathring\alpha,\mathbf B(\mathring\omega-\mathring\alpha)\rangle
=\Tr(\mathbf B\mathring\Lambda^{-1})+\langle\mathring\mu-\mathring\alpha,\mathbf B(\mathring\mu-\mathring\alpha)\rangle$;
taking the $R$-expectation of the difference of the two log densities gives \eqref{eq:link-gaussian-kl}. In
the mean-field case all operators are block diagonal, so traces and log-determinants are sums over blocks.
The general identity is \cref{thm:elbo-total-correlation-chain} with the product law $\nu^{\mathscr T}$,
and the total correlation of a Gaussian law is the sum of the summand entropies minus the joint entropy,
in which the factors $2\pi e$ cancel because the summand dimensions add up to $\dim\mathscr T$. (b) The
first claim is \cref{lem:link-growth}(a), and \cref{lem:link-growth}(c) gives
$\E_R\mathscr S\leq\mathrm c_1\E_Re^{\mathrm c_2\lVert\mathring\omega\rVert}<\infty$. With (a), the hypotheses of
the last sentence of \cref{prop:link-latent-normalization}(c) hold, and \eqref{eq:link-sector} is
$\log(dR/d\mathbb P^{\mathscr T}_\theta)=\log(dR/d\nu^{\mathscr T})+\mathscr S+\log Z_{\mathscr S}$ integrated
against $R$. (c) The baseline of the latent space is
$\mathbb P^{\mathscr T}_\theta(d\mathring\omega\given X)\mathbb P^{\mathsf Y}(dy\given\mathring\omega)$, and
\cref{lem:prob-kl-chain-rule}(a) writes its divergence from $Q$ as the first two terms of
\eqref{eq:link-gaussian-energy-form}, which are finite, so $Q$ is absolutely continuous with respect to
it; the negative part of $H_o$ is $Q$-integrable, and \eqref{eq:obs-collective-energy-form} of
\cref{thm:elbo-extended-gap}, applied on the latent space $\mathscr T\times\mathsf Y$, gives
\eqref{eq:link-gaussian-energy-form}. For a product law,
$\E_Q\lvert E_{a,o}\rvert\leq\mathrm c_1\E_Re^{\mathrm c_2\lVert\mathring\omega\rVert}\E_{Q^{\mathsf Y}}(1+\lVert y\rVert^2)<\infty$
by (b) and Tonelli's theorem. \hfill$\square$

The relative entropy \eqref{eq:link-sector} is the whole link sector: a loop energy placed in the prior
enters with its normalizer, as \cref{rem:gen-appended-energy} requires of any reciprocal energy that is to
be part of an evidence bound, and in the mean-field case the sector is a sum over edge copies of Gaussian
divergences from the anchors plus the recognition expectation of the loop action.

\section{The link field equation}
\label{sec:link-field}

Which condition does the recognition law of the generators satisfy where the collective free energy is
stationary? For a Gaussian recognition law and a given matter recognition law, the stress of the anchor
plus the expected gradient of the loop action equals the expected matter current, filtered through the
adjoint derivative of the exponential chart (\cref{prop:link-field-equation}); for a comparison record the
current is a precision-weighted prediction error times the source mean, minus the transported source
covariance (\cref{cor:link-delta-rule}). Without the generator prior the stationary problem degenerates,
at the sample level and at the law level (\cref{prop:link-law-level-energy,cor:link-prior-necessary}). The
recognition flow that relaxes to these conditions is \cref{prop:link-recognition-flow}. Matrices are real,
$\langle\mathbf S,\mathbf T\rangle=\Tr(\mathbf S^\top\mathbf T)$ is the trace pairing of square matrices of
one size, gradients of functions of matrices refer to it, and for a representation $\rho_x$ on a sample
fiber written in reference coordinates, $(\mathrm d\rho_x)^\top$ is the adjoint of $\mathrm d\rho_x$ for
$\langle\cdot,\cdot\rangle_{\mathsf G}$ on $\mathfrak g$ and the trace pairing.

\lemmaheading{Derivative of the exponential chart}{lem:link-exp-derivative}
For square matrices $\mathbf B,\mathbf T$ of one size put
\begin{equation}
\mathbf L(\mathbf B,\mathbf T)=\int_0^1e^{t\mathbf B}\mathbf Te^{(1-t)\mathbf B}dt.
\label{eq:link-exp-derivative}
\end{equation}
(a) For all $\mathbf B,\mathbf T$, in the operator norm,
$\lVert e^{\mathbf B+\mathbf T}-e^{\mathbf B}-\mathbf L(\mathbf B,\mathbf T)\rVert\leq\frac12\lVert\mathbf T\rVert^2e^{\lVert\mathbf B\rVert+\lVert\mathbf T\rVert}$,
so the matrix exponential is Fr\'echet differentiable at $\mathbf B$ with derivative
$\mathbf T\mapsto\mathbf L(\mathbf B,\mathbf T)$, and
$\lVert\mathbf L(\mathbf B,\mathbf T)\rVert\leq e^{\lVert\mathbf B\rVert}\lVert\mathbf T\rVert$. (b) For every
$\mathbf S$,
\begin{equation}
\langle\mathbf S,\mathbf L(\mathbf B,\mathbf T)\rangle=\langle\mathbf L(\mathbf B^\top,\mathbf S),\mathbf T\rangle.
\label{eq:link-exp-adjoint}
\end{equation}
(c) If a function $\mathbf h$ is differentiable at $e^{\mathbf B}$ with gradient $\nabla\mathbf h(e^{\mathbf B})$,
then $\mathbf B\mapsto\mathbf h(e^{\mathbf B})$ has gradient $\mathbf L(\mathbf B^\top,\nabla\mathbf h(e^{\mathbf B}))$
at $\mathbf B$, and $\mathring\xi\mapsto\mathbf h(\rho_x(\exp\mathring\xi))$ has gradient
$(\mathrm d\rho_x)^\top\mathbf L\bigl((\mathrm d\rho_x\mathring\xi)^\top,\nabla\mathbf h(\rho_x(\exp\mathring\xi))\bigr)$ at
$\mathring\xi$ for $\langle\cdot,\cdot\rangle_{\mathsf G}$. \status{ESTABLISHED}

\noindent\emph{Proof.}
For square matrices $\mathbf B$ and $\mathbf B'$, Duhamel's formula
$e^{\mathbf B'}-e^{\mathbf B}=\int_0^1e^{t\mathbf B'}(\mathbf B'-\mathbf B)e^{(1-t)\mathbf B}dt$ holds because the
integrand is the $t$-derivative of $e^{t\mathbf B'}e^{(1-t)\mathbf B}$, and with
$\lVert e^{t\mathbf B}\rVert\leq e^{t\lVert\mathbf B\rVert}$ it gives
$\lVert e^{\mathbf B'}-e^{\mathbf B}\rVert\leq\lVert\mathbf B'-\mathbf B\rVert e^{\max(\lVert\mathbf B'\rVert,\lVert\mathbf B\rVert)}$.
With $\mathbf B'=\mathbf B+\mathbf T$,
$e^{\mathbf B+\mathbf T}-e^{\mathbf B}-\mathbf L(\mathbf B,\mathbf T)=\int_0^1(e^{t(\mathbf B+\mathbf T)}-e^{t\mathbf B})\mathbf Te^{(1-t)\mathbf B}dt$,
and the same bound applied to $t(\mathbf B+\mathbf T)$ and $t\mathbf B$ bounds the integrand by
$t\lVert\mathbf T\rVert^2e^{t(\lVert\mathbf B\rVert+\lVert\mathbf T\rVert)+(1-t)\lVert\mathbf B\rVert}\leq t\lVert\mathbf T\rVert^2e^{\lVert\mathbf B\rVert+\lVert\mathbf T\rVert}$,
whose integral gives (a); the bound on $\mathbf L$ follows from
$\lVert e^{t\mathbf B}\rVert\lVert e^{(1-t)\mathbf B}\rVert\leq e^{\lVert\mathbf B\rVert}$.
For (b), cyclicity of the trace and $(e^{t\mathbf B})^\top=e^{t\mathbf B^\top}$ give
$\Tr(\mathbf S^\top e^{t\mathbf B}\mathbf Te^{(1-t)\mathbf B})=\Tr\bigl((e^{t\mathbf B^\top}\mathbf Se^{(1-t)\mathbf B^\top})^\top\mathbf T\bigr)$,
and integration over $t$ gives \eqref{eq:link-exp-adjoint}. Part (c) is the chain rule with (a) and (b);
for the representation, $\rho_x(\exp\mathring\xi)=\exp(\mathrm d\rho_x\mathring\xi)$ (\cref{lem:link-exp-naturality}),
its derivative in the direction $\mathring\eta$ is $\mathbf L(\mathrm d\rho_x\mathring\xi,\mathrm d\rho_x\mathring\eta)$, and
\eqref{eq:link-exp-adjoint}
moves $\mathbf L$ onto the gradient before the adjoint $(\mathrm d\rho_x)^\top$ is applied. \hfill$\square$

The integral form of the Fr\'echet derivative of the matrix exponential is classical \citep[Ch.~10]{Higham2008}.

\lemmaheading{Gaussian expectations and their derivatives}{lem:link-gaussian-gradients}
Let $\mathbf h$ be a twice continuously differentiable function on $\mathscr T$, or on one of its summands,
with
$\lvert\mathbf h(\mathring\omega)\rvert+\lVert\nabla\mathbf h(\mathring\omega)\rVert+\lVert\nabla^2\mathbf h(\mathring\omega)\rVert\leq\mathrm c_1e^{\mathrm c_2\lVert\mathring\omega\rVert}$.
Then $(\mu,\Sigma)\mapsto\E_{\mathcal N(\mu,\Sigma)}\mathbf h$ is continuously differentiable on the product
of the space and the open cone of positive-definite self-adjoint operators, and for every self-adjoint
$\mathring\Delta$
\begin{equation}
\nabla_\mu\E_{\mathcal N(\mu,\Sigma)}\mathbf h=\E_{\mathcal N(\mu,\Sigma)}\nabla\mathbf h,
\qquad
\frac{d}{dt}\E_{\mathcal N(\mu,\Sigma+t\mathring\Delta)}\mathbf h\Big|_{t=0}=\tfrac12\Tr\bigl(\mathring\Delta\E_{\mathcal N(\mu,\Sigma)}\nabla^2\mathbf h\bigr);
\label{eq:link-bonnet-price}
\end{equation}
that is, for the trace pairing of self-adjoint operators the gradient in $\Sigma$ is one half of
$\E_{\mathcal N(\mu,\Sigma)}\nabla^2\mathbf h$. \status{ESTABLISHED}

The proof is in \cref{sec:proofs-links}.

The two identities are the Gaussian derivative identities used for variational Gaussian approximations
in \citet[Eqs.~(18)--(19)]{OpperArchambeau2009}.

\propositionheading{Link field equation}{prop:link-field-equation}
Assume \cref{def:link-latent-model} with bounded record densities and $p(o\given\mathring\omega)>0$ for every
$\mathring\omega$, a $\sigma$-finite reference $\nu^{\mathsf Y}$ with respect to which every
$\mathbb P^{\mathsf Y}(\cdot\given\mathring\omega)$ has a density, and a matter recognition law $Q^{\mathsf Y}$
with density $q^{\mathsf Y}$ and $\E_{Q^{\mathsf Y}}\lvert\log q^{\mathsf Y}\rvert<\infty$. Let the matter energy
$\mathbf H_o(\mathring\omega,y)=-\log p_{\theta,V}(o,y\given X,\Theta(\mathring\omega))$, the negative logarithm
of the density \eqref{eq:gen-record-density} of the held-link joint at $\Theta(\mathring\omega)$, be twice
continuously differentiable in $\mathring\omega$ for every $y$ with
\begin{equation}
\E_{Q^{\mathsf Y}}\sup_{\lVert\mathring\omega\rVert\leq\mathring r}\Bigl(\lvert\mathbf H_o\rvert+\lVert\nabla_{\mathring\omega}\mathbf H_o\rVert
+\lVert\nabla^2_{\mathring\omega}\mathbf H_o\rVert\Bigr)(\mathring\omega,Y)\leq\mathrm c_1e^{\mathrm c_2\mathring r}
\qquad\text{for every }\mathring r\geq0.
\label{eq:link-matter-domination}
\end{equation}
Restrict recognition to the laws $R\otimes Q^{\mathsf Y}$ with $R=\mathcal N(\mathring\mu,\mathring\Lambda^{-1})$
Gaussian on $\mathscr T$, general or mean-field over the summands, and put
$\mathbf U(\mathring\omega)=\mathscr S(\mathring\omega)+\E_{Q^{\mathsf Y}}\mathbf H_o(\mathring\omega,Y)$. Then:

(a) $\mathbf U$ is twice continuously differentiable, with
$\lvert\mathbf U\rvert+\lVert\nabla\mathbf U\rVert+\lVert\nabla^2\mathbf U\rVert\leq\mathrm c_3e^{\mathrm c_4\lVert\mathring\omega\rVert}$,
and the free energy
\begin{equation}
\Fenergy_V[R\otimes Q^{\mathsf Y};X,o]=\KL(R\Vert\nu^{\mathscr T})+\E_R\mathbf U+\log Z_{\mathscr S}
+\E_{Q^{\mathsf Y}}\log q^{\mathsf Y}
\label{eq:link-field-free-energy}
\end{equation}
is finite and continuously differentiable in $(\mathring\mu,\mathring\Lambda)$.

(b) A pair $(\mathring\mu,\mathring\Lambda)$ is stationary for \eqref{eq:link-field-free-energy} exactly when
\begin{gather}
\mathring\Sigma^{-1}(\mathring\mu-\mathring\alpha)+\E_R\nabla\mathscr S
=-\E_{R\otimes Q^{\mathsf Y}}\nabla_{\mathring\omega}\mathbf H_o,
\label{eq:link-field-mean}\\
\mathring\Lambda=\mathring\Sigma^{-1}+\E_R\nabla^2\mathscr S+\E_{R\otimes Q^{\mathsf Y}}\nabla^2_{\mathring\omega}\mathbf H_o,
\label{eq:link-field-precision}
\end{gather}
and in the mean-field case exactly when the summand components of \eqref{eq:link-field-mean} and the
diagonal blocks of \eqref{eq:link-field-precision} hold. When the agent latent law does not read the links,
$\nabla_{\mathring\omega}\mathbf H_o=\sum_a\nabla_{\mathring\omega}E_{a,o}(X,\Theta(\mathring\omega),y_{\partial a})$, the sum over the
records whose kernels read a link, so the right side of \eqref{eq:link-field-mean} is minus the expected
gradient of the incident record energies.

(c) If $\mathbf H_o$ reads $\mathring\omega^x_e$ only through the represented twist
$\rho_x(\mathring\Theta^x_e)=\exp(\mathrm d\rho_x\mathring\omega^x_e)$, as a differentiable function of this matrix on
an open set of matrices that contains the represented twists, and its gradient
$\partial\mathbf H_o/\partial\rho_x(\mathring\Theta^x_e)$ is $Q^{\mathsf Y}$-integrable at every $\mathring\omega$, then
with the link current
\begin{equation}
\mathbf J^x_e(\mathring\omega)=-\E_{Q^{\mathsf Y}}\frac{\partial\mathbf H_o}{\partial\rho_x(\mathring\Theta^x_e)}(\mathring\omega,Y)
\label{eq:link-current}
\end{equation}
the $(x,e)$ component of \eqref{eq:link-field-mean} reads
\begin{equation}
(\mathring\Sigma^x_e)^{-1}(\mathring\mu^x_e-\mathring\alpha^x_e)+\E_R\nabla_{(x,e)}\mathscr S
=\E_R\Bigl[(\mathrm d\rho_x)^\top\mathbf L\bigl((\mathrm d\rho_x\mathring\omega^x_e)^\top,\mathbf J^x_e(\mathring\omega)\bigr)\Bigr]:
\label{eq:link-field-current}
\end{equation}
the expected stress of the generator prior equals the expected matter current transported by the adjoint
derivative of the exponential chart.

(d) Put
$\mathbf V(\mathring\omega)=\frac12\langle\mathring\omega-\mathring\alpha,\mathring\Sigma^{-1}(\mathring\omega-\mathring\alpha)\rangle+\mathbf U(\mathring\omega)$.
Then \eqref{eq:link-field-mean} is $\E_R\nabla\mathbf V=0$, and
$\E_{\mathcal N(\mathring\mu,\mathring\Lambda^{-1})}\nabla\mathbf V\to\nabla\mathbf V(\mathring\mu)$ as
$\mathring\Lambda^{-1}\to0$, uniformly for $\mathring\mu$ in compact sets. Every limit point $\mathring\mu_*$, as
$\mathring\Lambda^{-1}\to0$, of solutions of \eqref{eq:link-field-mean} satisfies $\nabla\mathbf V(\mathring\mu_*)=0$;
and if $\nabla\mathbf V(\mathring\mu_*)=0$ with $\nabla^2\mathbf V(\mathring\mu_*)$ invertible, then for every
positive-definite $\mathring\Lambda_0$ there are $t_0>0$ and a continuously differentiable curve
$t\mapsto\mathring\mu(t)$ on $(-t_0,t_0)$ with $\mathring\mu(0)=\mathring\mu_*$ such that
$(\mathring\mu(t),t^{-2}\mathring\Lambda_0)$ solves \eqref{eq:link-field-mean} for $0<\lvert t\rvert<t_0$. When
$\mathbf U$ is bounded below, the law
$R^\star(d\mathring\omega)\propto e^{-\mathbf U(\mathring\omega)}\nu^{\mathscr T}(d\mathring\omega)$ is the
unique minimizer of $\Fenergy_V[R\otimes Q^{\mathsf Y};X,o]$ over all laws $R$ on $\mathscr T$, and $-\mathbf V$ is
its log density up to a constant, so the limit condition $\nabla\mathbf V=0$ is the stationarity condition
of the most probable link given the matter recognition.

(e) No point mass is admissible: $\KL(\delta_{\mathring\mu}\Vert\nu^{\mathscr T})=+\infty$, so the free
energy of $\delta_{\mathring\mu}\otimes Q^{\mathsf Y}$ is $+\infty$, and \eqref{eq:link-field-free-energy} tends
to $+\infty$ as $\mathring\Lambda^{-1}\to0$ at a given $\mathring\mu$. \status{ESTABLISHED}

The proof is in \cref{sec:proofs-links}.

At a point law the matter current is evaluated at one link value and the condition $\nabla\mathbf V=0$ is the
most-probable-link equation; part (e) types it as a generative quantity at given matter recognition, in
the sense of \cref{prop:prob-law-point-point-mass}, and not as a recognition optimum.

\corollaryheading{The current of a comparison record}{cor:link-delta-rule}
In the setting of \cref{prop:link-field-equation}(c), let $x=b$, let $e:j\to i$ lie in $E_+$, and write samples
and links in the reference gauge, in which every frame is the identity, reached by the passive
re-expression by the inverse frames, which changes no twist and hence no generator
(\cref{prop:geo-edge-twist-invariance}). Suppose that the only term
of $\mathbf H_o$ that reads $\mathring\omega^b_e$ is either (a) a binary success record $a$ with scope $\{i,j\}$,
observed as a success, with
$E_{a,o}=\frac12(k_i-\rho_b(\mathring\Theta^b_e)k_j)^\top W_a(k_i-\rho_b(\mathring\Theta^b_e)k_j)$ and
$W_a\succ0$, or (b) the comparison record \eqref{eq:obs-selective-attention-record} of row $i$ at $o_i=0$ in
the per-receiver source-label model of \cref{def:obs-attention-label-model}, with query map $S^Q_i$, key map
$S^K_e=S^K_{\rm ref}\rho_b(\mathring\Theta^b_e)$ as in \cref{def:obs-selective-comparison}, tempered weight
$\widetilde W_i=W_i/\tau_i$, the source labels counted among the matter coordinates, and a matter recognition
law of the factorization \eqref{eq:obs-attention-recognition-factorization} whose row $\beta^Q_i$ does not
depend on $(\mathring\omega,y)$. Let the belief marginals of $Q^{\mathsf Y}$ have finite second moments, with
means $\mu_i,\mu_j$, covariance $\Sigma_j$ of $k_j$, and cross covariance
$\Sigma_{ij}=\E_{Q^{\mathsf Y}}[(k_i-\mu_i)(k_j-\mu_j)^\top]$, which vanishes for a product law. Then the
current \eqref{eq:link-current} is, in the two cases,
\begin{gather}
\mathbf J^b_e(\mathring\omega)
=W_a\bigl[(\mu_i-\rho_b(\mathring\Theta^b_e)\mu_j)\mu_j^\top+\Sigma_{ij}-\rho_b(\mathring\Theta^b_e)\Sigma_j\bigr],
\label{eq:link-delta-current}\\
\mathbf J^b_e(\mathring\omega)
=\beta^Q_{ij}(S^K_{\rm ref})^\top\widetilde W_i\bigl[(S^Q_i\mu_i-S^K_{\rm ref}\rho_b(\mathring\Theta^b_e)\mu_j)\mu_j^\top
+S^Q_i\Sigma_{ij}-S^K_{\rm ref}\rho_b(\mathring\Theta^b_e)\Sigma_j\bigr],
\label{eq:link-delta-current-attention}
\end{gather}
and \eqref{eq:link-field-current} holds for $e$ with this current. For a product law in case (a) the
current is the precision-weighted prediction error of the receiver times the source mean, minus the
precision-weighted transported source covariance. \status{ESTABLISHED}

\noindent\emph{Proof.}
For a matrix $\mathbf M$ in place of $\rho_b(\mathring\Theta^b_e)$, the gradient of
$\frac12(k_i-\mathbf Mk_j)^\top W_a(k_i-\mathbf Mk_j)$ in $\mathbf M$ for the trace pairing is
$-W_a(k_i-\mathbf Mk_j)k_j^\top$, which is $Q^{\mathsf Y}$-integrable by the second moments, and
$\E_{Q^{\mathsf Y}}[(k_i-\mathbf Mk_j)k_j^\top]=\mu_i\mu_j^\top+\Sigma_{ij}-\mathbf M(\Sigma_j+\mu_j\mu_j^\top)$, which
gives \eqref{eq:link-delta-current}. In case (b), \cref{prop:obs-attention-elbo}(b) with the latent state
$(\mathring\omega,y)$ shows that row $i$ contributes $\tau_i^{-1}\sum_j\beta^Q_{ij}\E D_{ij}$ together with terms
that do not read the links, because $\mathrm c_i$ of \cref{prop:obs-selective-receiver-update} is constant; the
source $j$ reached through $e$ has $D_{ij}=V_e$ of \cref{def:obs-selective-comparison}, whose gradient in
$\mathbf M$ is $-(S^K_{\rm ref})^\top W_i(S^Q_ik_i-S^K_{\rm ref}\mathbf Mk_j)k_j^\top$, and the same expectation gives
\eqref{eq:link-delta-current-attention}. The last sentence is
\cref{prop:link-field-equation}(c). \hfill$\square$

By \cref{cor:link-delta-rule}, the source of the link field equation \eqref{eq:link-field-current} has the
outer-product form of the delta rule that \citet[Sec.~4.2, Eq.~(24)]{Schlag2021} write for fast-weight
programmers, a line of work they trace in their Sec.~2 to \citet{Schmidhuber1992}: a weighted
prediction error of the receiver times the source mean. Three features have no counterpart there and are
consequences of the joint: the transported source covariance, the transport of the current by the
adjoint derivative of the exponential chart, and the average over the recognition law of the generators,
balanced by the stress of the generator prior. The correspondence is one of form between update rules; no
learning architecture is identified with the model.

\propositionheading{Law-level link energies and their stabilizers}{prop:link-law-level-energy}
Let $q_i,q_j$ be laws in the law fiber $\mathcal B_x$ of a channel, let
$D_e(\Theta)=\KL(q_i\Vert\widehat\rho_x(\Theta)q_j)$ for $\Theta\in G$ be the transported divergence
\eqref{eq:full-graph-transported-divergences} of a copy $e:j\to i$ as a function of its link value, and let
$G_q=\{g\in G:\widehat\rho_x(g)q=q\}$ be the stabilizer of $q$, a closed subgroup when
$g\mapsto\widehat\rho_x(g)q$ is continuous for the weak topology on laws, with Lie algebra $\mathfrak g_q$; under
\cref{hyp:pb-regular-models} it is the isotropy algebra of \eqref{eq:pb-orbit-fisher-form}.

(a) $D_e(h\Theta h')=D_e(\Theta)$ for all $\Theta\in G$, $h\in G_{q_i}$, $h'\in G_{q_j}$. Hence the set of
minimizers of $D_e$, and the set of its stationary points when $D_e$ is finite and continuously
differentiable on $G$, are unions of double cosets $G_{q_i}\Theta G_{q_j}$, and a minimizer or stationary
point is not isolated when $\dim\mathfrak g_{q_i}>0$ or $\dim\mathfrak g_{q_j}>0$. At every generator at which
the exponential map of $\mathfrak g$ is a local diffeomorphism, a minimizer or stationary point of the chart
energy $\mathring\omega^x_e\mapsto D_e(U_i\exp(\mathring\omega^x_e)U_j^{-1})$ is likewise not isolated when
$\dim\mathfrak g_{q_i}>0$ or $\dim\mathfrak g_{q_j}>0$.

(b) Let $G=\GL^+(K)$ act on $\R^K$ by its defining representation, and let
$q=\mathcal N(\mu,\Sigma)$ with $\Sigma\succ0$. Then
$G_q=\Sigma^{1/2}\{O\in\mathrm{SO}(K):O\Sigma^{-1/2}\mu=\Sigma^{-1/2}\mu\}\Sigma^{-1/2}$, of dimension
$(K-1)(K-2)/2$ if $\mu\neq0$ and $K(K-1)/2$ if $\mu=0$. The Mahalanobis norm
$\lvert\mu\rvert_\Sigma=(\mu^\top\Sigma^{-1}\mu)^{1/2}$ satisfies
$\lvert g\mu\rvert_{g\Sigma g^\top}=\lvert\mu\rvert_\Sigma$ for every $g$, and for $K\geq2$ a link $\Theta_0$ with
$\widehat\rho_x(\Theta_0)q_j=q_i$ exists exactly when $\lvert\mu_i\rvert_{\Sigma_i}=\lvert\mu_j\rvert_{\Sigma_j}$; the zero
set of $D_e$ is then the double coset $G_{q_i}\Theta_0G_{q_j}=\Theta_0G_{q_j}$. For $K\geq3$ every minimizer and
stationary point of $D_e$ is non-isolated, and for $K=2$ this holds when $\mu_i=0$ or $\mu_j=0$.

(c) If $\widehat\rho_x(\Theta_0)q_j=q_i$ and \cref{hyp:pb-regular-models} holds in the channel, then $D_e$ is
$C^3$ near $\Theta_0$ with zero value and differential at $\Theta_0$, and its Hessian at $\Theta_0$ on the
tangent vector of $t\mapsto\exp(t\mathring\xi)\Theta_0$ is the orbit Fisher form $\mathfrak k_{q_i}(\mathring\xi,\mathring\xi)$ of
\cref{cor:pb-orbit-fisher-form}; it is positive semidefinite with kernel the tangent space of the zero set.
For Gaussian laws in (b),
\begin{equation}
\mathfrak k_q(\mathring\xi,\mathring\xi)=\bigl\lvert\Sigma^{-1/2}\mathring\xi\mu\bigr\rvert^2
+\tfrac12\bigl\lVert\Sigma^{-1/2}\mathring\xi\Sigma^{1/2}+(\Sigma^{-1/2}\mathring\xi\Sigma^{1/2})^\top\bigr\rVert_{\mathrm F}^2,
\qquad\mathring\xi\in\mathfrak{gl}(K,\R).
\label{eq:link-gaussian-orbit-form}
\end{equation}
\status{ESTABLISHED}

The proof is in \cref{sec:proofs-links}.

\corollaryheading{Without the generator prior the link problem degenerates}{cor:link-prior-necessary}
(a) Let $G=\GL^+(K)$, $K\geq2$, act on the belief fiber by its defining representation, and in the reference
gauge let the record of \cref{cor:link-delta-rule}(a) be the only factor reading the link of $e:j\to i$.
For a matter recognition law with finite second moments, $\Sigma_j\succ0$, and $\Sigma_{ij}=0$, let the link be
moved as a parameter without a generator prior, with expected record energy
$\mathbf l(\mathbf M)=\E_{Q^{\mathsf Y}}\frac12(k_i-\mathbf Mk_j)^\top W_a(k_i-\mathbf Mk_j)$ for real
$K\times K$ matrices $\mathbf M$. Then:
(i) $\mathbf l$ is a strictly convex quadratic whose unique stationary point and minimizer is
\begin{equation}
\mathbf M^\star=\mu_i\mu_j^\top(\Sigma_j+\mu_j\mu_j^\top)^{-1}
=\frac{\mu_i\mu_j^\top\Sigma_j^{-1}}{1+\mu_j^\top\Sigma_j^{-1}\mu_j},
\label{eq:link-rank-one}
\end{equation}
of rank at most one and therefore singular; with a cross covariance it is instead the second-moment
regression $\E[k_ik_j^\top]\E[k_jk_j^\top]^{-1}$, which can be invertible. (ii) $\mathbf l$ has no stationary
point on $\GL(K)$; its infimum over $\GL^+(K)$ equals $\mathbf l(\mathbf M^\star)$, is not attained, and every
minimizing sequence in $\GL^+(K)$ converges to $\mathbf M^\star$. (iii) A generator $\mathring\omega^b_e$ is
stationary for $\mathring\omega^b_e\mapsto\mathbf l(\exp\mathring\omega^b_e)$ only where
$\mathbf L(\mathring\omega^b_e,\cdot)$ is singular, and such generators exist: for $K=2$,
$\mu_i=\mu_j=\begin{psmallmatrix}1\\0\end{psmallmatrix}$, $\Sigma_j=W_a=I_2$, the generator
$\mathring\omega^b_e=-(\log3)I_2+2\pi\begin{psmallmatrix}0&-1\\1&0\end{psmallmatrix}$ is stationary, with link
$\exp\mathring\omega^b_e=\frac13I_2\neq\mathbf M^\star$. (iv) If $\mu_j^\top\Sigma_j^{-1}\mu_i\geq0$, then
$\mathbf M^\star$ is a limit of chart links $\exp\mathring\omega^b_e$, so the infimum over chart links equals
$\mathbf l(\mathbf M^\star)$, is not attained, and is approached only along links converging to
$\mathbf M^\star$; if $\mu_j^\top\Sigma_j^{-1}\mu_i<0$, no sequence of chart links converges to $\mathbf M^\star$, the
infimum over chart links exceeds $\mathbf l(\mathbf M^\star)$, and for $K=2$ it is not attained.

(b) If the free energy depends on a link value only through a law-point energy $D_e$ of
\cref{prop:link-law-level-energy}, finite and continuously differentiable on $G$, and the link is moved as a
parameter without a generator prior, its stationary link values form a union of double cosets, none
isolated when $\dim\mathfrak g_{q_i}+\dim\mathfrak g_{q_j}>0$; for nondegenerate Gaussian laws under $\GL^+(K)$,
where $D_e$ is finite and smooth, this holds for every pair of laws when $K\geq3$.

(c) With the generator prior: (i) under the hypotheses of \cref{prop:link-field-equation} with $\mathbf U$
bounded below, $\mathbf V$ attains its minimum on $\mathscr T$, and
$(\mathring\mu,\mathring\Lambda)\mapsto\Fenergy_V[\mathcal N(\mathring\mu,\mathring\Lambda^{-1})\otimes Q^{\mathsf Y};X,o]$ has
compact sublevel sets in the product of $\mathscr T$ and the open cone of positive-definite precisions, or of
block-diagonal ones in the mean-field case, and attains its minimum there; (ii) if $\mathbb P^{\mathsf Y}$ is a
nondegenerate Gaussian law on a real vector space $\mathsf Y$ of finite dimension that does not read the links,
and every record energy is nonnegative, lower semicontinuous in $(\mathring\omega,y)$, and bounded by
$\mathrm c_1(1+\lVert y\rVert^2)e^{\mathrm c_2\lVert\mathring\omega\rVert}$, then the free energy restricted to the product
Gaussian laws $\mathcal N(\mathring\mu,\mathring\Lambda^{-1})\otimes\mathcal N(\mu^{\mathsf Y},\Sigma^{\mathsf Y})$ has compact sublevel
sets in $(\mathring\mu,\mathring\Lambda,\mu^{\mathsf Y},\Sigma^{\mathsf Y})$ and attains its minimum.
\status{ESTABLISHED}

The proof is in \cref{sec:proofs-links}.

At the sample level the degeneracy of (a) is tied to a product matter recognition, since a correlated one
can make the stationary matrix invertible, while the coercivity of (c) comes from the anchor term of the
prior at every width $\mathring\Sigma$. Part (a)(iii) shows that the chart itself creates stationary generators
at its folds: a gradient flow in generator coordinates started at such a generator does not move, although
its link value is not stationary for the link energy.

\section{Held links as a limit, and conditional statements}
\label{sec:link-held}

How does the theory with latent generators relate to the theory in which links are structural data or
parameters? As the prior covariance shrinks at a held anchor, the infimum of the latent-link free energy
converges to the infimum of the held-link free energy at the anchor links, and near-minimizing link
recognition laws concentrate at the anchor (\cref{thm:link-held-limit}); at every width, the conditional
laws given the generators are held-link laws, so statements proved for every link value transfer
conditionally, while statements about mixtures over the generators need not
(\cref{prop:link-conditional-transfer}).

\theoremheading{Held links are the zero-width limit}{thm:link-held-limit}
Hold $\theta$ except for the prior covariances, which are $\mathring\Sigma=\mathring\sigma^2\mathring\Sigma_0$ for
$\mathring\sigma>0$ and a block-diagonal $\mathring\Sigma_0\succ0$, and hold $X$, $o$ and the anchor $\mathring\alpha$.
Write $\nu^{\mathscr T}_{\mathring\sigma}$, $\mathbb P^{\mathscr T}_{\mathring\sigma}$, $Z_{\mathscr S,\mathring\sigma}$ and
$\Fenergy_{\mathring\sigma}$ for the base law, the generator prior, its normalizer, and the free energy of the
latent-link model at width $\mathring\sigma$. Let $\mathcal Q'$ be a nonempty class of laws on $\mathsf Y$, put
$\mathbf F_Q(\mathring\omega)=\Fenergy_V[Q;X,\Theta(\mathring\omega),o]$ for $Q\in\mathcal Q'$ and
$\mathbf F(\mathring\omega)=\inf_{Q\in\mathcal Q'}\mathbf F_Q(\mathring\omega)$, and assume:
(i) the record densities are bounded, $L_o\leq\mathrm c_{\mathfrak A}$, and $p(o\given\mathring\omega)>0$ for every
$\mathring\omega$;
(ii) $\mathbf F$ is finite and lower semicontinuous at $\mathring\alpha$;
(iii) a subclass $\mathcal Q'_0\subseteq\mathcal Q'$ with $\inf_{Q\in\mathcal Q'_0}\mathbf F_Q(\mathring\alpha)=\mathbf F(\mathring\alpha)$
consists of laws for which $\mathbf F_Q$ is continuous at $\mathring\alpha$ and bounded by
$\mathrm c_Qe^{\mathrm c_Q\lVert\mathring\omega\rVert}$.
Then:

(a) As $\mathring\sigma\to0$, the infimum of $\Fenergy_{\mathring\sigma}$ over the product laws $R\otimes Q$ with
$Q\in\mathcal Q'$ and $R$ in any class of laws on $\mathscr T$ that contains $\nu^{\mathscr T}_{\mathring\sigma}$, and its
infimum over all laws $R(d\mathring\omega)Q_{\mathring\omega}(dy)$ whose kernel takes values in $\mathcal Q'$, both
converge to
\begin{equation}
\mathbf F(\mathring\alpha)=\inf_{Q\in\mathcal Q'}\Fenergy_V\bigl[Q;X,\Theta(\mathring\alpha),o\bigr],
\qquad
\Theta^x_e(\mathring\alpha)=U_i\exp(\mathring\alpha^x_e)U_j^{-1},
\label{eq:link-held-limit}
\end{equation}
the infimum over $\mathcal Q'$ of the held-link free energy at the anchor links.

(b) Under (i) alone, every family of laws $R_{\mathring\sigma}(d\mathring\omega)Q_{\mathring\sigma,\mathring\omega}(dy)$ with
$\sup_{\mathring\sigma}\Fenergy_{\mathring\sigma}<\infty$ satisfies
$R_{\mathring\sigma}\{\lVert\mathring\omega-\mathring\alpha\rVert\geq\mathring\epsilon\}\to0$ for every $\mathring\epsilon>0$; under (i)
and (ii),
every such family with kernels in $\mathcal Q'$ and $\Fenergy_{\mathring\sigma}\to\mathbf F(\mathring\alpha)$ satisfies in
addition $\KL(R_{\mathring\sigma}\Vert\mathbb P^{\mathscr T}_{\mathring\sigma})\to0$.

(c) Under (i) alone, if $p(o\given\cdot)$ is continuous at $\mathring\alpha$, then the evidence of the
latent-link model converges to $p_{\theta,V}(o\given X,\Theta(\mathring\alpha))$, and the link marginal
$\boldsymbol\Pi^{\mathscr T}_{\mathring\sigma}$ of its posterior satisfies
$\boldsymbol\Pi^{\mathscr T}_{\mathring\sigma}\{\lVert\mathring\omega-\mathring\alpha\rVert\geq\mathring\epsilon\}\to0$ for every
$\mathring\epsilon>0$.

(d) Under (i), hypothesis (ii) holds when $\mathcal Q'=\{Q_\lambda:\lambda\in\mathbb U\}$ for a subset $\mathbb U$ of a
Euclidean space, $(\mathring\omega,\lambda)\mapsto\mathbf F_{Q_\lambda}(\mathring\omega)$ is lower semicontinuous,
$\mathbf F(\mathring\alpha)<\infty$, and for some $\mathring\epsilon_0>0$ and some $\mathrm c>\mathbf F(\mathring\alpha)$ the set
of $\lambda$ with $\mathbf F_{Q_\lambda}(\mathring\omega)\leq\mathrm c$ for some
$\lVert\mathring\omega-\mathring\alpha\rVert<\mathring\epsilon_0$ has compact closure in $\mathbb U$. \status{ESTABLISHED}

The proof is in \cref{sec:proofs-links}.

The present held-link theory is therefore the zero-width limit of the latent one in value: the prior and
the recognition law narrow together, so the link sector costs nothing in the limit, in contrast with a
recognition law that narrows under a prior of constant width, whose cost diverges
(\cref{prop:link-field-equation}(e)).

\propositionheading{Conditional statements transfer; link-marginal statements need not}{prop:link-conditional-transfer}
Assume the hypotheses of \cref{prop:link-latent-normalization}(c).
(a) For every $\mathring\omega$ the conditional law of the posterior given $\mathring\omega$ is the held-link
posterior $\boldsymbol\Pi_{\theta,V,o,X,\Theta(\mathring\omega)}$, and for every law
$Q=R(d\mathring\omega)Q_{\mathring\omega}(dy)$ the conditional free energy
$-\log p(o\given\mathring\omega)+\KL(Q_{\mathring\omega}\Vert\boldsymbol\Pi_{\mathring\omega})$ is the held-link free energy
$\Fenergy_V[Q_{\mathring\omega};X,\Theta(\mathring\omega),o]$. If $Q'=R(d\mathring\omega)Q'_{\mathring\omega}(dy)$ has the same link
marginal and both free energies are finite, then
\begin{equation}
\Fenergy_V[Q';X,o]-\Fenergy_V[Q;X,o]
=\int\Bigl(\Fenergy_V\bigl[Q'_{\mathring\omega};X,\Theta(\mathring\omega),o\bigr]-\Fenergy_V\bigl[Q_{\mathring\omega};X,\Theta(\mathring\omega),o\bigr]\Bigr)R(d\mathring\omega).
\label{eq:link-transfer-potential}
\end{equation}
(b) Let $\mathbb O$ be a measurable set of link families, and consider a property of a pair consisting of
held-link structural data and a law on $\mathsf Y$, expressed through the held-link joint, evidence,
posterior and free energy. If the property holds for $((X,\Theta),Q'')$ for every $\Theta\in\mathbb O$ and
every $Q''$ in a class $\mathcal Q''$, then it holds for $((X,\Theta(\mathring\omega)),Q_{\mathring\omega})$ at every
$\mathring\omega$ with $\Theta(\mathring\omega)\in\mathbb O$ and $Q_{\mathring\omega}\in\mathcal Q''$, hence for $R$-almost
every $\mathring\omega\in\Theta^{-1}(\mathbb O)$ when the kernel lies in $\mathcal Q''$ $R$-almost surely. For instance, for
$G=\GL^+(K)$ in its defining representation and a closed walk $p$ in which the copies of some reversal pair
occur with nonzero net count, the generators at which $\lvert\det\rho_x(\mathrm{Hol}^x(p))\rvert=1$ form a
hyperplane of $\mathscr T$, so for every $R$ absolutely continuous with respect to Lebesgue measure the
conclusion of \cref{cor:cg-holonomy-noncompact-gaussian} for a block containing $p$ holds conditionally on
$\mathring\omega$ for $R$-almost every $\mathring\omega$.
(c) A property of laws on $\mathsf Y$ transfers from the conditional laws to the link-marginal laws, namely
the matter marginal $\int Q_{\mathring\omega}R(d\mathring\omega)$ of a recognition law, the matter marginal
$\int\boldsymbol\Pi_{\mathring\omega}\boldsymbol\Pi^{\mathscr T}(d\mathring\omega)$ of the posterior, and their images under a coarse
channel that does not read the links, when the set of laws with the property is preserved by mixtures, and
need not otherwise. For $G=\GL^+(1)$ acting on $\R$ by multiplication, two agents with
$k_j\sim\mathcal N(0,1)$ and $k_i\given k_j\sim\mathcal N(\rho_b(\Theta^b_e)k_j,1)$ for the copy $e:j\to i$, trivial
frames, $\mathscr S=0$, $\nu^{\mathscr T}=\mathcal N(0,1)$, and no record, every held-link posterior is Gaussian,
while the $k_i$-marginal of the latent-link posterior has
$\E k_i^4=3(1+2e^2+e^8)>3(1+e^2)^2=3(\E k_i^2)^2$, so it is not Gaussian. \status{ESTABLISHED}

\noindent\emph{Proof.}
(a) The first claim is the disintegration of \cref{prop:link-latent-normalization}(c), and the second is
\cref{def:elbo-extended} at $(X,\Theta(\mathring\omega))$. Finite free energies make both integrals in
\eqref{eq:link-decomposition} finite, because the integrands are bounded below and the divergence term is
common, so subtraction gives \eqref{eq:link-transfer-potential}. (b) By (a) the conditional objects at
$\mathring\omega$ are the held-link objects at $(X,\Theta(\mathring\omega))$, and kernels are determined up to
$R$-null sets. For the example, $\det e^{\mathbf B}=e^{\operatorname{tr}\mathbf B}$ (triangularize $\mathbf B$ over
$\mathbb C$), so by \eqref{eq:geo-edge-twist-holonomy}
$\det\rho_x(\mathrm{Hol}^x(p))$ is the exponential of the linear functional
$\mathring\omega\mapsto\sum_{e\in E_+}\operatorname{tr}\mathring\omega^x_e$ weighted by the net count of each $e$ in $p$, which is
nonzero by assumption, and the set where the determinant has modulus one is the kernel of this functional,
which is Lebesgue null. Off it the holonomy has $\lvert\det\rvert\neq1$, and
\cref{cor:cg-holonomy-noncompact-gaussian} applies at every such link value. (c) The marginals are mixtures
over $\mathring\omega$ of the conditional laws, and a link-independent kernel commutes with mixing, which gives
the positive clause. In the example there are no records, so the posterior is the prior; given the
generator, $(k_i,k_j)$ is Gaussian with $\E[k_i^2\given\mathring\omega]=1+e^{2\mathring\omega}$, and a centered Gaussian
has $\E k^4=3(\E k^2)^2$, so $\E k_i^4=3\E(1+e^{2\mathring\omega})^2=3(1+2e^2+e^8)$ and $\E k_i^2=1+e^2$, by
$\E e^{2\mathring\omega}=e^2$ and $\E e^{4\mathring\omega}=e^8$; strict inequality is strict Jensen, since
$e^{2\mathring\omega}$ is not almost surely constant. \hfill$\square$

A coarse channel that reads the links, such as the tree transports of
\cref{def:full-graph-coarse-links}, is a kernel on $(\mathring\omega,y)$ and is covered by (b) conditionally on
$\mathring\omega$; its link-marginal image is again a mixture.

\section{Directed and undirected matter}
\label{sec:link-induced}

Does matter generate a link action of its own? In the record extension it does not: the generator marginal
is the generator prior, and matter reaches the links only through the evidence. When generators and matter
are normalized jointly as an undirected model, the generator marginal acquires minus the logarithm of the
matter normalizer (\cref{prop:link-induced-action}).

\propositionheading{Induced link action}{prop:link-induced-action}
(a) For the joint \eqref{eq:link-latent-joint},
$\int_{\mathsf Y\times\mathsf O}\mathbb P_{\theta,V}(d\mathring\omega,dy,do\given X)=\mathbb P^{\mathscr T}_\theta(d\mathring\omega\given X)$,
and with record densities the link marginal of the posterior is proportional to
$p(o\given\mathring\omega)\mathbb P^{\mathscr T}_\theta(d\mathring\omega\given X)$.
(b) Let the potentials $\psi_i(\cdot;X,\Theta)$, $\psi_{ij}(\cdot;X,\Theta)$, $\psi_a(\cdot;X,\Theta)$ of
\cref{def:gen-undirected-model} depend jointly measurably on the link family, write
$\widetilde p_{X,\Theta}$ and $Z_{X,\Theta}$ for \eqref{eq:gen-gibbs-potential} and \eqref{eq:gen-gibbs-partition}
at $(X,\Theta)$, and normalize generators and matter jointly: with
\begin{equation}
\mathfrak Z=\int_{\mathscr T}e^{-\mathscr S(\mathring\omega)}Z_{X,\Theta(\mathring\omega)}\nu^{\mathscr T}(d\mathring\omega),
\label{eq:link-joint-undirected-normalizer}
\end{equation}
assume $0<\mathfrak Z<\infty$, and take the law with density
$\mathfrak Z^{-1}e^{-\mathscr S(\mathring\omega)}\widetilde p_{X,\Theta(\mathring\omega)}(o,y)$ with respect to
$\nu^{\mathscr T}\otimes\nu^{\mathsf Y}\otimes\nu^{\mathsf O}$. It is a probability law,
$Z_{X,\Theta(\mathring\omega)}<\infty$ for $\nu^{\mathscr T}$-almost every $\mathring\omega$, its generator marginal is
\begin{equation}
\mathfrak Z^{-1}e^{-\mathscr S(\mathring\omega)-\mathscr S^{\mathrm{ind}}(\mathring\omega)}\nu^{\mathscr T}(d\mathring\omega),
\qquad
\mathscr S^{\mathrm{ind}}(\mathring\omega)=-\log Z_{X,\Theta(\mathring\omega)},
\label{eq:link-induced-marginal}
\end{equation}
with $\mathscr S^{\mathrm{ind}}(\mathring\omega)\in(-\infty,+\infty]$ for $\nu^{\mathscr T}$-almost every
$\mathring\omega$, and on $\{0<Z_{X,\Theta(\mathring\omega)}<\infty\}$ its conditional law given $\mathring\omega$ is the undirected model
\eqref{eq:gen-gibbs-density} at $(X,\Theta(\mathring\omega))$. A record extension with densities is the case
$Z_{X,\Theta}=1$ for every $\Theta$, in which $\mathscr S^{\mathrm{ind}}=0$. The condition $\mathfrak Z<\infty$ is not
implied by $Z_{\mathscr S}<\infty$: for $G=\GL^+(1)$, trivial frames, $\mathscr S=0$, two scalar agents with
standard Gaussian self potentials, and the edge potential $\exp(\rho_b(\Theta_e)k_1k_2)$,
$Z_{X,\Theta}=(1-\rho_b(\Theta_e)^2)^{-1/2}$ when $\rho_b(\Theta_e)<1$ and $Z_{X,\Theta}=+\infty$ when
$\rho_b(\Theta_e)=e^{\mathring\omega_e}\geq1$, a set of positive $\nu^{\mathscr T}$-measure, so $\mathfrak Z=+\infty$.
\status{ESTABLISHED}

\noindent\emph{Proof.}
(a) By Tonelli's theorem, integrate the records, each a kernel of unit mass, and then $y$ against the
probability kernel $\mathbb P^{\mathsf Y}(\cdot\given\mathring\omega)$; the posterior marginal is that of
\cref{prop:link-latent-normalization}(c). (b) Tonelli's theorem gives
$\int e^{-\mathscr S}\widetilde p_{X,\Theta(\mathring\omega)}d(\nu^{\mathsf Y}\otimes\nu^{\mathsf O})=e^{-\mathscr S(\mathring\omega)}Z_{X,\Theta(\mathring\omega)}$,
whose $\nu^{\mathscr T}$-integral is $\mathfrak Z$, so the density integrates to one; finiteness of $\mathfrak Z$ with
$e^{-\mathscr S}>0$ forces $Z_{X,\Theta(\mathring\omega)}<\infty$ almost everywhere, the generator marginal is
$\mathfrak Z^{-1}e^{-\mathscr S}Z_{X,\Theta(\mathring\omega)}\nu^{\mathscr T}$, which is \eqref{eq:link-induced-marginal} with the
value $+\infty$ of $\mathscr S^{\mathrm{ind}}$ where $Z_{X,\Theta(\mathring\omega)}=0$, and dividing the joint density by
the marginal density gives \eqref{eq:gen-gibbs-density}. The densities of a record extension are
conditional densities of unit mass in $(y,o)$ for every link family (\cref{prop:gen-exact-normalization}),
so as potentials they have $Z_{X,\Theta}=1$. In the example,
$\iint(2\pi)^{-1}e^{-(k_1^2+k_2^2)/2+\rho_b(\Theta_e)k_1k_2}dk_1dk_2=\det\begin{psmallmatrix}1&-\rho_b(\Theta_e)\\-\rho_b(\Theta_e)&1\end{psmallmatrix}^{-1/2}$
for $\rho_b(\Theta_e)<1$, the integral diverges for $\rho_b(\Theta_e)\geq1$ by
\cref{prop:gen-gibbs-counterexample}, and the Gaussian law $\nu^{\mathscr T}$ on $\mathfrak g=\R$ charges
$\{\mathring\omega_e\geq0\}$. \hfill$\square$

For Gaussian self potentials and pairwise potentials of the connection-Laplacian form
\eqref{eq:cg-connection-laplacian-energy}, the matter normalizer is a Gaussian determinant, so the induced
action is one half the log-determinant of the connection Laplacian plus the self term, and its expansion in
closed walks produces loop products of the links; Supplement G develops both.

\section{Inference dynamics of the links}
\label{sec:link-inference-dynamics}

The generators are latent coordinates, so they move by the dynamics of every recognition coordinate
of the program, descent of the collective free energy. This section writes that descent for a
Gaussian recognition law of the generators jointly with the agents' recognition coordinates, shows
that it is a natural-gradient history whose explicit Euler step is the Gaussian Bayesian learning
rule (\cref{prop:link-recognition-flow}), gives the exact M-coordinate of the anchor
(\cref{prop:link-anchor-m-coordinate}), and reads the supplied link rate of
\cref{prop:obs-slow-model-link-dichotomy} in this setting (\cref{rem:link-rates}). In this section
and the next two, the links of one channel are latent, or both channels share them under the
common-link specialization \eqref{eq:geo-common-link}, and the channel index is dropped. The link
space $\mathscr T=\mathfrak g^{E_+}$ carries the product inner product
$\langle\mathring\omega,\mathring\omega'\rangle=\sum_{e\in E_+}\langle\mathring\omega_e,\mathring\omega'_e\rangle_{\mathsf G}$
of \cref{def:link-generator-prior}, its norm $\lVert\cdot\rVert$, and the Lebesgue measure of that definition;
self-adjointness, traces, and determinants of operators on $\mathscr T$ refer to this inner product.
Where a representation $\rho$ enters, $\mathrm d\rho$ is injective and $\mathfrak g$ carries the
pulled-back inner product of \cref{def:link-loop-defect}, so that
$\lVert\mathring\omega_e\rVert_{\mathsf G}=\lVert\mathrm d\rho\mathring\omega_e\rVert_{\mathsf G}$. For a
matrix $\mathbf H$ on the represented fiber,
$\mathbf H^*=(\mathsf G^{\mathrm{fib}})^{-1}\mathbf H^\top\mathsf G^{\mathrm{fib}}$ is its adjoint for the
fiber metric, so that $\langle\mathbf H,\mathbf H'\rangle_{\mathsf G}=\Tr(\mathbf H^*\mathbf H')$ in
\eqref{eq:link-trace-form} and left or right multiplication by $\mathbf H$ has as adjoint the left or
right multiplication by $\mathbf H^*$; $\operatorname{pr}_{\mathfrak g}$ is the
$\langle\cdot,\cdot\rangle_{\mathsf G}$-orthogonal projection of matrices onto
$\mathrm d\rho(\mathfrak g)$, read in $\mathfrak g$.

\propositionheading{Natural-gradient recognition flow of the generators}{prop:link-recognition-flow}
In the latent-link model of \cref{def:link-latent-model}, with bounded record densities and
$p(o\given\mathring\omega)>0$ for every $\mathring\omega$, let the recognition law be
$R\otimes Q_\lambda$. Here $R=\mathcal N(\mathring\mu,\mathring\Lambda^{-1})$ is a Gaussian law on
$\mathscr T$ with mean $\mathring\mu\in\mathscr T$ and precision $\mathring\Lambda$ in the open cone
$\mathring\Xi_{\succ0}$ of positive-definite elements of the space $\mathring\Xi$ of self-adjoint
operators that preserve each summand of a given orthogonal decomposition of $\mathscr T$: one summand
admits every positive-definite precision, the summands $\mathfrak g$ of the edge copies give the
products $\bigotimes_{e\in E_+}\mathcal N(\mathring\mu_e,\mathring\Lambda_e^{-1})$, and the lines of an
orthonormal basis give the diagonal precisions. The law $Q_\lambda$ is the agents' recognition law,
with coordinates $\lambda$ in a collective configuration manifold $\mathcal Q_V$
(\cref{def:hist-configuration-manifold}) that satisfies \cref{hyp:hist-standing-configuration-tier}
with metric $\mathsf G_V$. By \eqref{eq:link-decomposition} and
\cref{prop:link-gaussian-recognition-vfe}(b) with $Q_{\mathring\omega}=Q_\lambda$ for every
$\mathring\omega$, the collective free energy at $R\otimes Q_\lambda$ is
\begin{equation}
F(\mathring\mu,\mathring\Lambda,\lambda)
=\E_R\bigl[\mathfrak i_\lambda\bigr]+\tfrac12\log\det\mathring\Lambda+\log Z_{\mathscr S}
-\tfrac12\dim\mathscr T\log(2\pi e),
\label{eq:link-gaussian-free-energy}
\end{equation}
with the generator integrand
\begin{equation}
\mathfrak i_\lambda(\mathring\omega)
=\Fenergy_V\bigl[Q_\lambda;X,\Theta(\mathring\omega),o\bigr]+\mathscr S(\mathring\omega)
+\tfrac12\bigl\langle\mathring\omega-\mathring\alpha,\mathring\Sigma^{-1}(\mathring\omega-\mathring\alpha)\bigr\rangle
+\tfrac12\log\det\bigl(2\pi\mathring\Sigma\bigr).
\label{eq:link-generator-integrand}
\end{equation}
Assume that $F$ is $C^2$ on $\mathscr T\times\mathring\Xi_{\succ0}\times\mathcal Q_V$ and that each
$\mathfrak i_\lambda$ is $C^2$ with
$\lvert\mathfrak i_\lambda\rvert+\lVert\nabla\mathfrak i_\lambda\rVert+\lVert\nabla^2\mathfrak i_\lambda\rVert\leq\mathrm c_1e^{\mathrm c_2\lVert\mathring\omega\rVert}$
for constants $\mathrm c_1,\mathrm c_2$. Write $\mathring\mu\otimes\mathring\mu$ for the operator
$\mathring\omega\mapsto\langle\mathring\mu,\mathring\omega\rangle\mathring\mu$ and
$\operatorname{pr}_{\mathring\Xi}$ for the orthogonal projection onto $\mathring\Xi$ in the trace inner
product.
\begin{enumerate}
\item[(a)] The laws $\mathcal N(\mathring\mu,\mathring\Lambda^{-1})$ with
$\mathring\Lambda\in\mathring\Xi_{\succ0}$ form a minimal regular exponential family with natural
parameter $\mathring\vartheta=(\mathring\Lambda\mathring\mu,-\tfrac12\mathring\Lambda)$, expectation
parameter
$\mathring\varpi=(\mathring\mu,\mathring\Lambda^{-1}+\operatorname{pr}_{\mathring\Xi}(\mathring\mu\otimes\mathring\mu))$,
and log-partition function $\mathring\Gamma$, hence a regular statistical manifold. Its Fisher metric
is smooth and positive definite, and in the coordinates $(\mathring\mu,\mathring\Lambda)$ it is
\begin{equation}
g^F\bigl((\delta\mathring\mu,\delta\mathring\Lambda),(\delta\mathring\mu,\delta\mathring\Lambda)\bigr)
=\langle\delta\mathring\mu,\mathring\Lambda\delta\mathring\mu\rangle
+\tfrac12\Tr\bigl(\mathring\Lambda^{-1}\delta\mathring\Lambda\mathring\Lambda^{-1}\delta\mathring\Lambda\bigr).
\label{eq:link-gaussian-fisher}
\end{equation}
The product $\mathcal Q$ of this family with $\mathcal Q_V$, with the metric
$\mathsf G=g^F\oplus\mathsf G_V$, satisfies \cref{hyp:hist-standing-configuration-tier}.
\item[(b)] The natural-gradient field $-\operatorname{grad}^{\mathsf G}F$ of
\eqref{eq:hist-fisher-gradient-definition} on $\mathcal Q$ is $C^1$, so its maximal integral curves
exist and are unique. Along each of them $\tfrac{d}{dt}F=-\lVert\operatorname{grad}^{\mathsf G}F\rVert_{\mathsf G}^2\leq0$,
and each one through a noncritical point has nowhere-vanishing velocity and satisfies
\eqref{eq:hist-vfe-ray-representative} with $\mathsf a\equiv1$, so it represents a natural-gradient
history (\cref{def:hist-oriented-history}).
\item[(c)] The flow is of first order. Its velocity is a function of the configuration alone and
vanishes exactly at the critical points of $F$. It obeys the force balance
$\mathsf G(\dot Q,\cdot)=-dF$, in which $\mathsf G$ is the friction tensor, and
$\tfrac{d}{dt}F=-\mathsf G(\dot Q,\dot Q)$. The free energy strictly decreases along every
nonconstant solution, so no nonconstant solution is periodic. At a critical point $Q_\star$ with
invertible Hessian the linearized field is $-\mathsf G(Q_\star)^{-1}\nabla^2F(Q_\star)$, which is
diagonalizable with real eigenvalues, so no linearized mode oscillates.
\item[(d)] In the chart $(\mathring\vartheta,\lambda)$, the generator block of the explicit Euler step
of the flow with step $\Delta t>0$ is
\begin{equation}
\mathring\Lambda'=(1-\Delta t)\mathring\Lambda+\Delta t\operatorname{pr}_{\mathring\Xi}\E_R\bigl[\nabla^2\mathfrak i_\lambda\bigr],
\qquad
\mathring\mu'=\mathring\mu-\Delta t\mathring\Lambda'^{-1}\E_R\bigl[\nabla\mathfrak i_\lambda\bigr],
\label{eq:link-blr}
\end{equation}
whenever $\mathring\Lambda'$ is positive definite, for instance when $0<\Delta t<1$ and
$\operatorname{pr}_{\mathring\Xi}\E_R[\nabla^2\mathfrak i_\lambda]$ is positive semidefinite. For one summand
this is the Gaussian form of the Bayesian learning rule \citep[Eq.~(12)]{KhanRue2023} applied to
$\mathfrak i_\lambda$, and for the lines of an orthonormal basis it is its diagonal-precision form, which
\citet[Sec.~4.2, Eq.~(37)]{KhanRue2023} state with the expected gradient and Hessian replaced by their
values at the mean.
\item[(e)] Let $R'=\mathcal N(\mathring\mu',\mathring\Lambda'^{-1})$ with expectation parameter
$\mathring\varpi'$, and let $\widetilde F(\mathring\varpi)$ denote $F$ at fixed $\lambda$ in the
expectation chart. If $\mathring\Lambda'$ is positive definite and
$\nabla^2\widetilde F\preceq\mathrm c_{\mathsf G}\nabla^2\mathring\Gamma^\star$ on the segment from
$\mathring\varpi$ to $\mathring\varpi'$ for some $\mathrm c_{\mathsf G}>0$, where
$\mathring\Gamma^\star$ is the convex conjugate of $\mathring\Gamma$ and $\nabla^2\mathring\Gamma^\star$ is
the Fisher metric in the expectation chart, then
\begin{equation}
F(R')\leq F(R)-\Delta t^{-1}\KL(R\Vert R')-\bigl(\Delta t^{-1}-\mathrm c_{\mathsf G}\bigr)\KL(R'\Vert R),
\label{eq:link-blr-descent}
\end{equation}
so a step with $0<\Delta t\leq1/\mathrm c_{\mathsf G}$ decreases $F$ unless $R'=R$. For one summand
and a quadratic integrand
$\mathfrak i_\lambda(\mathring\omega)=\tfrac12\langle\mathring\omega,\mathbf Q\mathring\omega\rangle-\langle\mathbf b,\mathring\omega\rangle+\mathrm c_3$
with $\mathbf Q$ positive definite, the hypothesis holds with $\mathrm c_{\mathsf G}=1$ everywhere, and
the step $\Delta t=1$ gives $R'=\mathcal N(\mathbf Q^{-1}\mathbf b,\mathbf Q^{-1})$, the minimizer of $F$
over Gaussian $R$.
\end{enumerate}
\status{ESTABLISHED}

The proof is in \cref{sec:proofs-links}.

For links, the latent option of \cref{def:obs-selective-interface-learning} asks of a random relation
coordinate a prior that is a probability law, dynamics in the augmented state, and a regular configuration
metric or a justified quotient, and the present construction discharges these requirements. The generator
prior of \cref{def:link-generator-prior}, in the joint of \cref{def:link-latent-model}, supplies the first for
every Lie group $G$, compact or not (\cref{prop:link-generator-prior,prop:link-latent-normalization}), with
directed matter, and with undirected matter normalized jointly whenever the normalizer
\eqref{eq:link-joint-undirected-normalizer} is finite (\cref{prop:link-induced-action}); the dynamics and the
metric are those of \cref{prop:link-recognition-flow}.

\propositionheading{Exact M-coordinate of the anchor}{prop:link-anchor-m-coordinate}
Let the generator prior of \cref{def:link-generator-prior} have base law
$\nu^{\mathscr T}_{\mathring\alpha}=\bigotimes_{e\in E_+}\mathcal N(\mathring\alpha_e,\mathring\Sigma_e)$
and an action $\mathscr S\geq0$ that does not depend on $\mathring\alpha$, so that the prior is
$\mathbb P^{\mathscr T}_{\mathring\alpha}=Z_{\mathscr S}(\mathring\alpha)^{-1}e^{-\mathscr S}\nu^{\mathscr T}_{\mathring\alpha}$
with $Z_{\mathscr S}(\mathring\alpha)=\int e^{-\mathscr S}d\nu^{\mathscr T}_{\mathring\alpha}$, and let
no other factor of the joint depend on $\mathring\alpha$. Hold the recognition coupling
$\mathbb Q_{V,o,X}$ fixed, let its generator marginal $R$ have finite second moments, and assume the
finiteness hypotheses of \cref{prop:elbo-exact-m-coordinate} for every anchor.
\begin{enumerate}
\item[(a)] Up to terms that do not depend on $\mathring\alpha$, the objective of
\eqref{eq:elbo-m-coordinate} is
\begin{equation}
\Lelbo^{\mathrm{anc}}(\mathring\alpha)=-\tfrac12\E_R\bigl\langle\mathring\omega-\mathring\alpha,\mathring\Sigma^{-1}(\mathring\omega-\mathring\alpha)\bigr\rangle-\log Z_{\mathscr S}(\mathring\alpha),
\label{eq:link-anchor-objective}
\end{equation}
which contains the normalizer of the prior, and
\begin{equation}
\nabla_{\mathring\alpha_e}\Lelbo^{\mathrm{anc}}=\mathring\Sigma_e^{-1}\bigl(\E_R[\mathring\omega_e]-\E_{\mathbb P^{\mathscr T}_{\mathring\alpha}}[\mathring\omega_e]\bigr).
\label{eq:link-anchor-gradient}
\end{equation}
\item[(b)] The function $\Lelbo^{\mathrm{anc}}$ is strictly concave and has exactly one maximizer
$\mathring\alpha^{\mathrm{new}}$, the unique anchor at which the prior mean equals the recognition
mean, $\E_{\mathbb P^{\mathscr T}_{\mathring\alpha^{\mathrm{new}}}}[\mathring\omega]=\E_R[\mathring\omega]$. By
\cref{prop:elbo-exact-m-coordinate} it is the exact M-coordinate on the anchor at the given
covariances.
\item[(c)] If $\mathscr S=0$, as when every loop coefficient $\mathrm c_p$ vanishes, then
$\mathring\alpha^{\mathrm{new}}=\E_R[\mathring\omega]$, which is $\mathring\mu$ for a Gaussian
generator recognition law. If
$\mathscr S(\mathring\omega)=\tfrac12\sum_p\mathrm c_p\lVert\mathbf C_p\mathring\omega\rVert_{\mathsf G}^2$,
where $\mathbf C_p\mathring\omega=(\mathbf C\mathring\omega)_p$ is the loop sum of the cycle map
\eqref{eq:link-cycle-map} and $\mathbf C_p^\top$ is the adjoint of $\mathbf C_p:\mathscr T\to\mathfrak g$,
as for the loop defects of the translation instance in its affine representation
(\cref{prop:link-translation-instance}), then
\begin{equation}
\mathring\alpha^{\mathrm{new}}=\Bigl(I+\mathring\Sigma\sum_p\mathrm c_p\mathbf C_p^\top\mathbf C_p\Bigr)\E_R[\mathring\omega],
\label{eq:link-anchor-quadratic}
\end{equation}
which differs from $\E_R[\mathring\omega]$ exactly when $\mathbf C_p\E_R[\mathring\omega]\neq0$ for
some loop $p$ with $\mathrm c_p>0$.
\item[(d)] If the covariances $\mathring\Sigma_e$ range as well over the positive-definite
self-adjoint operators on $\mathfrak g$, the objective is strictly concave in the coordinates
$(\mathring\Sigma_e^{-1}\mathring\alpha_e,-\tfrac12\mathring\Sigma_e^{-1})_{e\in E_+}$, it has at most
one maximizer, and a maximizer matches the mean and the covariance of every edge generator under
the prior and under $R$. If $\mathscr S=0$ and every $\operatorname{Cov}_R(\mathring\omega_e)$ is
positive definite, the maximizer exists and is
$(\E_R[\mathring\omega_e],\operatorname{Cov}_R(\mathring\omega_e))_{e\in E_+}$. With a loop term it
need not exist: in the translation instance with $\mathsf G^{\mathrm{fib}}=I_2$, on one agent $i$
with one reversal pair of edge copies $e,\bar e:i\to i$, the loop $p=(e)$ with
$\mathrm c_p=\mathrm c>0$, and $R=\mathcal N(\mathring\mu,\mathrm v)$ with $\mathrm v\geq1/\mathrm c$,
the objective has no maximizer over $(\mathring\alpha_e,\mathring\Sigma_e)\in\R\times(0,\infty)$, and
its supremum is approached as $\mathring\Sigma_e\to\infty$.
\end{enumerate}
\status{ESTABLISHED}

The proof is in \cref{sec:proofs-links}.

\remarkheading{The supplied link rate in the latent reading}{rem:link-rates}
Read \cref{def:rg-formation-cycle-experiment} with every frame at the reference point, so that each
twist equals its link (\cref{prop:geo-edge-twist-invariance}).
(i) Its link translations, written $\check\omega_i$ in \cref{prop:obs-slow-model-link-dichotomy}, are
generators of the translation instance (\cref{prop:link-translation-instance}): $(\R,+)$ acting on the
scalar fibers by translation, in its affine representation with $\mathsf G^{\mathrm{fib}}=I_2$, for
which $\mathfrak g=\R$, $\langle\mathrm x,\mathrm y\rangle_{\mathsf G}=\mathrm x\mathrm y$, and every loop
defect is the quadratic $\tfrac12(\sum_a\check\omega_{e_a})^2$. The link block $\mathcal N(0,I_4)$ of the
configuration prior is the generator prior of \cref{def:link-generator-prior} with anchor $0$,
$\mathring\Sigma_e=1$, and no loop term; the loop energy of the experiment enters through its
records. For a law of the configuration, the generator
recognition law $R$ is its $\check\omega$-marginal, and its mean is the generator recognition mean
$\mathring\mu$.
(ii) Moving the link translations at the supplied rate $\varepsilon_{\check\omega}$ is therefore, on
the held matching sheet and for $\varepsilon_{\mathrm m},\varepsilon_{\check\omega}>0$, the
rate-separated flow \eqref{eq:obs-rate-separated-flow} with the generator recognition means in the
slow block, $s=(\bar m,\mathring\mu)$ and
$\mathsf G_{\mathrm s}=I_4\oplus(\varepsilon_{\mathrm m}/\varepsilon_{\check\omega})I_4$, by
\cref{prop:obs-slow-model-link-dichotomy}(i).
(iii) The anchor $\mathring\alpha$ is a coordinate of the generative parameter $\theta$, so it
belongs to the slow block $s=(s_{\mathrm m},\theta)$ of \cref{def:obs-fast-slow-blocks}: it moves the
joint between evaluations, by accepted M-coordinates (\cref{def:elbo-generalized-m-acceptance}), of
which \cref{prop:link-anchor-m-coordinate} gives the exact one, or by the rate-separated flow, and
never inside an evaluation (\cref{req:gen-typing-prohibition}). In the experiment no generative
parameter is learned and the anchor stays at $0$.
(iv) The recognition flow of \cref{prop:link-recognition-flow} does not supply the separation of
(ii). In the notation of \cref{def:rg-formation-cycle-experiment}, take the held sheet $S^{(0)}$, the
Gaussian configuration laws
$\mathcal N(\mu,\Sigma)\otimes\delta_{S^{(0)}}$ on $\R^{12}$ with their Fisher metric, of the form
\eqref{eq:link-gaussian-fisher}, and the free energy $\Fenergy_{\mathrm T}/\mathrm T$ of
\cref{prop:rg-formation-endogenous-cycle}, at the temperature $\mathrm T$ of that proposition. Then
$\nabla^2U_{S^{(0)}}\succeq\mathrm TI$, the natural-gradient flow leaves the slice
$\Sigma=\mathrm T(\nabla^2U_{S^{(0)}})^{-1}$ invariant, and on it $\mu(t)=e^{-t}\mu(0)$: belief means,
model means, and link means relax at one common rate. From
$\mu(0)=\mathbf x^\circ$, every residual is $r^b_i(t)=r^m_i(t)=-e^{-t}$, so at no time is
$\lvert r^b_i\rvert\leq2\varepsilon_{\mathrm m}<\tfrac12\leq\lvert r^m_i\rvert$, and the
belief-agreeing, model-disagreeing interval of \cref{prop:obs-slow-model-link-dichotomy}(ii) does not
occur. \status{ESTABLISHED}

The proof is in \cref{sec:proofs-links}.

On the slice of \cref{rem:link-rates}(iv) the Fisher geometry of the recognition family assigns every
coordinate the same rate, so there the separation in \cref{prop:obs-slow-model-link-dichotomy} comes
from its supplied mobilities alone, and a slow link block is a declaration of the kind that
\cref{def:obs-fast-slow-blocks} allows.

\section{Link histories across learning episodes}
\label{sec:link-histories}

Across learning episodes the generators form a history, and the latent route of
\cref{def:full-graph-episode-joint} gives that history one joint law. This section puts a Gibbs prior
on generator histories whose kinetic term is the loop defect of a time-like loop
(\cref{def:link-history-prior}), proves that it is a latent route with the generators in the parent
readout (\cref{prop:link-history-prior}), and shows that its increment precision is the inertial mass
of the most probable history (\cref{prop:link-map-inertia}). Throughout,
$\mathfrak Q(\mathbf H)=\tfrac12\lVert\mathbf H-I\rVert_{\mathsf G}^2$ for a matrix $\mathbf H$ on the
represented fiber, so that the loop defect \eqref{eq:link-loop-defect} of a closed walk $p$ is
$\mathfrak Q_p=\mathfrak Q(\rho(\mathring\Theta_p))$.

\definitionheading{Generator history prior}{def:link-history-prior}
Index learning episodes by $n=0,\ldots,n_1$ as in \cref{def:full-graph-episode-joint}, let
$\mathring\omega^{(n)}=(\mathring\omega^{(n)}_e)_{e\in E_+}\in\mathscr T$ be the generators of episode
$n$, with $\mathring\omega^{(n)}_{\bar e}=-\mathring\omega^{(n)}_e$, and write
$\mathring\omega^{(0:n_1)}$ for a history. Fix the generator prior of \cref{def:link-generator-prior},
with base law $\nu^{\mathscr T}=\bigotimes_{e\in E_+}\mathcal N(\mathring\alpha_e,\mathring\Sigma_e)$ and
action $\mathscr S=\sum_p\mathrm c_p\mathfrak Q_p$ (the statements below use only that $\mathscr S$ is
finite, measurable, and nonnegative), and a mass $\mathbf m>0$. Frames are static at $c_*$: the copies
of agent $i$ at episodes $n$ and $n+1$ are joined by the identity link. For an edge copy $e:j\to i$,
the time-like loop of $e$ between episodes $n$ and $n+1$ is the closed walk that runs from $j$ at
episode $n+1$ to $j$ at episode $n$, along $e$ at episode $n$, from $i$ at episode $n$ to $i$ at
episode $n+1$, and along $\bar e$ at episode $n+1$. Its reference holonomy
\eqref{eq:link-reference-holonomy}, with the identity twist for each time-like step, is
$\mathring\Theta^{(n+1\leftarrow n)}_e=\exp(-\mathring\omega^{(n+1)}_e)\exp(\mathring\omega^{(n)}_e)$, with
the represented holonomy
\begin{equation}
\rho\bigl(\mathring\Theta^{(n+1\leftarrow n)}_e\bigr)=e^{-\mathrm d\rho\mathring\omega^{(n+1)}_e}e^{\mathrm d\rho\mathring\omega^{(n)}_e}.
\label{eq:link-timelike-holonomy}
\end{equation}
The kinetic term of $e$ between episodes $n$ and $n+1$ is either the time-like loop defect
$\mathfrak U^{(n)}_e=\mathfrak Q(\rho(\mathring\Theta^{(n+1\leftarrow n)}_e))$ or the Gaussian increment
$\mathfrak U^{(n)}_e=\tfrac12\lVert\mathring\omega^{(n+1)}_e-\mathring\omega^{(n)}_e\rVert_{\mathsf G}^2$,
and $\mathfrak U(\mathring\omega^{(n+1)},\mathring\omega^{(n)})=\sum_{e\in E_+}\mathfrak U^{(n)}_e$. The
generator history prior is the law on $\mathscr T^{n_1+1}$
\begin{equation}
\mathbb P^{\mathrm{hist}}\bigl(d\mathring\omega^{(0:n_1)}\bigr)
=Z_{\mathrm{hist}}^{-1}\exp\Bigl[-\sum_{n=0}^{n_1}\mathscr S\bigl(\mathring\omega^{(n)}\bigr)
-\mathbf m\sum_{n=0}^{n_1-1}\mathfrak U\bigl(\mathring\omega^{(n+1)},\mathring\omega^{(n)}\bigr)\Bigr]
\prod_{n=0}^{n_1}\nu^{\mathscr T}\bigl(d\mathring\omega^{(n)}\bigr),
\label{eq:link-history-prior}
\end{equation}
with $Z_{\mathrm{hist}}$ the integral of the exponential against $\bigotimes_n\nu^{\mathscr T}$. The
anchor, the prior covariances, the loop coefficients, and the mass are coordinates of the generative
parameter $\theta$.

\propositionheading{The generator history prior}{prop:link-history-prior}
In \cref{def:link-history-prior}, with either kinetic term:
\begin{enumerate}
\item[(a)] $0<Z_{\mathrm{hist}}\leq Z_{\mathscr S}^{n_1+1}\leq1$, so
$\mathbb P^{\mathrm{hist}}$ is a probability law, absolutely continuous with respect to
$\bigotimes_n\nu^{\mathscr T}$ with density at most $Z_{\mathrm{hist}}^{-1}$; every function
integrable under $\bigotimes_n\nu^{\mathscr T}$, in particular every function bounded by
$\mathrm c_1e^{\mathrm c_2\lVert\mathring\omega^{(0:n_1)}\rVert}$, is integrable under it.
\item[(b)] With the backward functions $\mathring Z^{\leftarrow}_{n_1}\equiv1$ and
\begin{equation}
\mathring Z^{\leftarrow}_n(\mathring\omega)=\int\exp\bigl[-\mathscr S(\mathring\omega')-\mathbf m\mathfrak U(\mathring\omega',\mathring\omega)\bigr]\mathring Z^{\leftarrow}_{n+1}(\mathring\omega')\nu^{\mathscr T}(d\mathring\omega'),
\qquad n<n_1,
\label{eq:link-backward-functions}
\end{equation}
which are measurable with values in $(0,1]$, $\mathbb P^{\mathrm{hist}}$ is the law of the Markov chain
with root and transition kernels
\begin{equation}
\begin{gathered}
\mathbb P^{(0)}_{\mathrm{hist}}(d\mathring\omega)=Z_{\mathrm{hist}}^{-1}e^{-\mathscr S(\mathring\omega)}\mathring Z^{\leftarrow}_0(\mathring\omega)\nu^{\mathscr T}(d\mathring\omega),\\
\mathbb P^{(n+1\leftarrow n)}_{\mathrm{hist}}(d\mathring\omega'\given\mathring\omega)
=\frac{e^{-\mathscr S(\mathring\omega')-\mathbf m\mathfrak U(\mathring\omega',\mathring\omega)}\mathring Z^{\leftarrow}_{n+1}(\mathring\omega')}{\mathring Z^{\leftarrow}_n(\mathring\omega)}\nu^{\mathscr T}(d\mathring\omega').
\end{gathered}
\label{eq:link-history-kernels}
\end{equation}
These kernels are the conditional laws of the Gibbs law \eqref{eq:link-history-prior} and depend on
$n$ and on the horizon $n_1$ through $\mathring Z^{\leftarrow}_{n+1}$. The chain whose kernels normalize each
increment factor separately, started from the law proportional to $e^{-\mathscr S}\nu^{\mathscr T}$,
has the law \eqref{eq:link-history-prior} exactly when $\mathring Z^{\leftarrow}_{n_1-1}$ is constant
$\nu^{\mathscr T}$-almost everywhere, and this fails for one scalar generator with
$\nu^{\mathscr T}=\mathcal N(0,1)$, $\mathscr S=0$, the Gaussian increment, and $n_1=1$. With tower states
$\mathbf Y^{(n)}=(\mathring\omega^{(n)},y^{(n)})$ at depth zero, parent readout
$\chi^{\mathrm{ep}}(\mathbf Y^{(n)})=\mathring\omega^{(n)}$, the agents' latent law of each episode given
by a Markov kernel $\mathbb P^{\mathrm{ag}}_n(dy^{(n)}\given\mathring\omega^{(n)},X)$, and records of
episode $n$ that read $\mathbf Y^{(n)}$ only, the joint is a latent route of
\cref{def:full-graph-episode-joint}, with root
$\mathbb P^{(0)}_{\mathrm{hist}}(d\mathring\omega^{(0)})\mathbb P^{\mathrm{ag}}_0(dy^{(0)}\given\mathring\omega^{(0)},X)$ and transitions
$\mathbb P^{(n\leftarrow n-1)}_{\mathrm{hist}}(d\mathring\omega^{(n)}\given\mathring\omega^{(n-1)})\mathbb P^{\mathrm{ag}}_n(dy^{(n)}\given\mathring\omega^{(n)},X)$,
each a composition of kernels. Hence
\cref{prop:full-graph-episode-joint} applies: the joint has unit mass and is acyclic with forward
cross-episode edges, the root is required, the predictive prior of the generators at episode $n+1$ is
the filtered law of $\mathring\omega^{(n)}$ pushed through $\mathbb P^{(n+1\leftarrow n)}_{\mathrm{hist}}$, and, on the
finite-density domain with observation densities, the free energy decomposes along $n$ as in
\eqref{eq:full-graph-episode-decomposition}.
\item[(c)] The time-like links carry no variable. The twist of the identity link of agent $i$ between
two episodes is $e_G$, so $\mathbb P^{\mathrm{hist}}$ is a law on the spatial generator histories
alone, and each kinetic term is a function of the generators of one edge at two consecutive episodes.
In lattice gauge theory this configuration is the temporal gauge, in which a time-like plaquette
couples two space-like links at consecutive times \citep[Ch.~15, Eq.~(15.9)]{Creutz1983}.
\end{enumerate}
\status{ESTABLISHED}

The proof is in \cref{sec:proofs-links}.

\propositionheading{Inertia of the most probable history}{prop:link-map-inertia}
Let $\mathfrak V_n:\mathscr T\to\R$, $n=0,\ldots,n_1$, be differentiable, let $\mathbf m>0$ and
$\mathfrak U$ be either kinetic term of \cref{def:link-history-prior}, and let
\begin{equation}
\mathscr S^{\mathrm{hist}}\bigl(\mathring\omega^{(0:n_1)}\bigr)=\sum_{n=0}^{n_1}\mathfrak V_n\bigl(\mathring\omega^{(n)}\bigr)
+\mathbf m\sum_{n=0}^{n_1-1}\mathfrak U\bigl(\mathring\omega^{(n+1)},\mathring\omega^{(n)}\bigr).
\label{eq:link-history-action}
\end{equation}
For the history prior, $\mathfrak V_n=\mathscr S+\tfrac12\langle\cdot-\mathring\alpha,\mathring\Sigma^{-1}(\cdot-\mathring\alpha)\rangle$
makes $\mathscr S^{\mathrm{hist}}$ the negative log density of \eqref{eq:link-history-prior} relative to Lebesgue measure up
to a constant; adding to $\mathfrak V_n$ a matter term of episode $n$, such as the negative log
evidence of its records given the generators or their expected energy under a matter recognition law,
gives the corresponding posterior or conditional objective. A most probable history is a minimizer of
$\mathscr S^{\mathrm{hist}}$, with free ends or with ends held at given values.
\begin{enumerate}
\item[(a)] For the Gaussian increment, every critical point satisfies
\begin{equation}
\mathbf m\bigl(2\mathring\omega^{(n)}-\mathring\omega^{(n+1)}-\mathring\omega^{(n-1)}\bigr)+\nabla\mathfrak V_n\bigl(\mathring\omega^{(n)}\bigr)=0,
\qquad 0<n<n_1,
\label{eq:link-map-inertia}
\end{equation}
and, with free ends,
$\mathbf m(\mathring\omega^{(0)}-\mathring\omega^{(1)})+\nabla\mathfrak V_0(\mathring\omega^{(0)})=0$ and
$\mathbf m(\mathring\omega^{(n_1)}-\mathring\omega^{(n_1-1)})+\nabla\mathfrak V_{n_1}(\mathring\omega^{(n_1)})=0$.
With unit step in $n$, \eqref{eq:link-map-inertia} is the central-difference form of
$\mathbf m\ddot{\mathring\omega}=\nabla\mathfrak V$, Newton's equation with mass $\mathbf m$ in the
potential $-\mathfrak V$, the Euclidean sign. If each $\mathfrak V_n$ is continuous and
$\sum_n\mathfrak V_n(\mathring\omega^{(n)})\to\infty$ as
$\lVert\mathring\omega^{(0:n_1)}\rVert\to\infty$, as when each $\mathfrak V_n$ is at least the anchor
quadratic minus a constant, a most probable history exists.
\item[(b)] For the time-like loop defect, write
$\mathring\Theta^{\leftarrow}_e=\exp(-\mathring\omega'_e)\exp(\mathring\omega_e)$ for a later family
$\mathring\omega'$ and an earlier family $\mathring\omega$, so that $\rho(\mathring\Theta^{\leftarrow}_e)$ is
\eqref{eq:link-timelike-holonomy} with $\mathring\omega'$ and $\mathring\omega$ in place of
$\mathring\omega^{(n+1)}$ and $\mathring\omega^{(n)}$, and let $\mathbf L$ be the Fr\'echet derivative
\eqref{eq:link-exp-derivative}. The gradients of $\mathfrak U(\mathring\omega',\mathring\omega)$ in its
earlier and later arguments are
\begin{equation}
\begin{aligned}
\bigl(\partial_{\mathring\omega}\mathfrak U\bigr)_e&=\operatorname{pr}_{\mathfrak g}\mathbf L\bigl((\mathrm d\rho\mathring\omega_e)^*,(e^{-\mathrm d\rho\mathring\omega'_e})^*\bigl(\rho(\mathring\Theta^{\leftarrow}_e)-I\bigr)\bigr),\\
\bigl(\partial_{\mathring\omega'}\mathfrak U\bigr)_e&=-\operatorname{pr}_{\mathfrak g}\mathbf L\bigl(-(\mathrm d\rho\mathring\omega'_e)^*,\bigl(\rho(\mathring\Theta^{\leftarrow}_e)-I\bigr)(e^{\mathrm d\rho\mathring\omega_e})^*\bigr),
\end{aligned}
\label{eq:link-timelike-gradients}
\end{equation}
and every critical point satisfies, for $0<n<n_1$,
\begin{equation}
\mathbf m\bigl[\partial_{\mathring\omega}\mathfrak U\bigl(\mathring\omega^{(n+1)},\mathring\omega^{(n)}\bigr)+\partial_{\mathring\omega'}\mathfrak U\bigl(\mathring\omega^{(n)},\mathring\omega^{(n-1)}\bigr)\bigr]+\nabla\mathfrak V_n\bigl(\mathring\omega^{(n)}\bigr)=0.
\label{eq:link-map-inertia-defect}
\end{equation}
For $\mathrm r>0$ let $\mathrm c_{\mathrm r}=\sqrt2\sup\lVert D^3\mathfrak U(\mathring\omega',\mathring\omega)\rVert$ over
$\lVert\mathring\omega\rVert,\lVert\mathring\omega'\rVert\leq\mathrm r$. On that set
$\lVert\partial_{\mathring\omega}\mathfrak U-(\mathring\omega-\mathring\omega')\rVert$ and
$\lVert\partial_{\mathring\omega'}\mathfrak U-(\mathring\omega'-\mathring\omega)\rVert$ are at most
$\mathrm c_{\mathrm r}\mathrm r\lVert\mathring\omega'-\mathring\omega\rVert$, so at a history in the ball of radius $\mathrm r$
the left side of \eqref{eq:link-map-inertia-defect} differs from that of \eqref{eq:link-map-inertia}
by at most
$\mathbf m\mathrm c_{\mathrm r}\mathrm r(\lVert\mathring\omega^{(n+1)}-\mathring\omega^{(n)}\rVert+\lVert\mathring\omega^{(n)}-\mathring\omega^{(n-1)}\rVert)$.
At a given $\mathring\omega$ and small increments $\delta\mathring\omega$,
$\mathfrak U(\mathring\omega+\delta\mathring\omega,\mathring\omega)=\tfrac12\sum_e\lVert\mathbf L(-\mathrm d\rho\mathring\omega_e,\mathrm d\rho\delta\mathring\omega_e)e^{\mathrm d\rho\mathring\omega_e}\rVert_{\mathsf G}^2+O(\lVert\delta\mathring\omega\rVert^3)$,
so the time-like defect gives the increment the mass tensor
$\delta\mathring\omega\mapsto\mathbf m\sum_e\lVert\mathbf L(-\mathrm d\rho\mathring\omega_e,\mathrm d\rho\delta\mathring\omega_e)e^{\mathrm d\rho\mathring\omega_e}\rVert_{\mathsf G}^2$,
which depends on the generators and reduces to $\mathbf m\lVert\delta\mathring\omega\rVert^2$ at
$\mathring\omega=0$.
\item[(c)] Let $\mathscr T$ be one-dimensional, as for one reversal pair of edge copies of the
translation instance, identified with $\R$ isometrically, take the Gaussian increment, hold the ends
at $\mathring\omega^{(0)}=\mathrm a<\mathrm b=\mathring\omega^{(n_1)}$, and let $\mathfrak V_n=\mathfrak V$
for $0<n<n_1$ be continuous with $\mathfrak V\geq\mathfrak V(\mathrm a)$ on $(-\infty,\mathrm a]$ and
$\mathfrak V\geq\mathfrak V(\mathrm b)$ on $[\mathrm b,\infty)$, as for a double well with minima at
$\mathrm a$ and $\mathrm b$. Then a most probable history exists, and every most probable history is
nondecreasing in $n$; if moreover $\mathfrak V$ is differentiable with
$\mathfrak V'(\mathrm a)=\mathfrak V'(\mathrm b)=0$, every most probable history is strictly
increasing. The most probable path between the wells is a monotone kink and does not oscillate.
\end{enumerate}
\status{ESTABLISHED}

The proof is in \cref{sec:proofs-links}.

The inertia of \cref{prop:link-map-inertia} belongs to the smoothed history, a joint optimization over
$n$, while the recognition flow of \cref{prop:link-recognition-flow} is of first order in its
parameter. Recognition dynamics with velocity variables arise when the generative model itself
declares the motion of its states, as in the generalized coordinates of motion of
\citet[p.~852, Eq.~(13)]{Friston2008DEM}, whose ensemble dynamics couple each state to its motion;
the history prior declares no such variables. The Gaussian increment term of
\eqref{eq:link-history-action} is, up to a constant, the negative log transition density of a random
walk with increment covariance $\mathbf m^{-1}I$, the discrete-time form of the path weight of a
driftless diffusion.

\section{The linearized link equation}
\label{sec:link-linear}

Near zero generators the stationarity condition of the generator history becomes a linear equation
whose operator collects every term of the construction: the anchor, the loops, the episodes, and the
matter. Its positivity makes the linearized history unique, and the neglected terms are of second
order in the sources (\cref{prop:link-linearized-equation}).

\propositionheading{The linearized link equation}{prop:link-linearized-equation}
In the latent-link model of \cref{def:link-latent-model}, let the generators carry the history prior
of \cref{def:link-history-prior} with zero anchor, prior covariance $\mathring\Sigma$, action
$\mathscr S=\sum_p\mathrm c_p\mathfrak Q_p$, mass $\mathbf m$, and either kinetic term, and write samples
in the reference gauge of \cref{cor:link-delta-rule}. At each episode $n$ let $\mathfrak A_n$ be a
finite set of binary success records (\cref{hyp:local-interaction-kernels}), observed as successes;
record $a$ has a scope $\{i,j\}$ joined by an edge copy $e:j\to i$ of $E_+$ or of its reversal, with
$\mathring\omega_{\bar e}=-\mathring\omega_e$, and the energy of \cref{cor:link-delta-rule}(a),
\begin{equation}
E_{a,o}(k;\mathring\omega)=\tfrac12\bigl(k_i-e^{\mathrm d\rho\mathring\omega_e}k_j\bigr)^\top W_a\bigl(k_i-e^{\mathrm d\rho\mathring\omega_e}k_j\bigr),
\qquad W_a\succ0,
\label{eq:link-comparison-record}
\end{equation}
where $i$, $j$, and $e$ are those of $a$. Record $a$ enters the free energy of its episode with a
weight $\beta_a\in[0,1]$ that does not depend on the generators: $\beta_a=1$ for a record that is
always read, and a weight below one arises when a record enters through a recognized source label,
as in \cref{cor:link-delta-rule}(b). Hold a product Gaussian matter recognition under which $k_i$ at
episode $n$ has mean $\mu^{(n)}_i$ and covariance $\Sigma^{(n)}_i$, and let $\mathscr S^{\mathrm{hist}}$ be
the history action \eqref{eq:link-history-action} with
$\mathfrak V_n=\tfrac12\langle\cdot,\mathring\Sigma^{-1}\cdot\rangle+\mathscr S+\sum_{a\in\mathfrak A_n}\beta_a\E[E_{a,o}(k;\cdot)]$,
that is,
\begin{equation}
\begin{aligned}
\mathscr S^{\mathrm{hist}}(\mathring\omega)={}&\sum_{n=0}^{n_1}\Bigl[\tfrac12\bigl\langle\mathring\omega^{(n)},\mathring\Sigma^{-1}\mathring\omega^{(n)}\bigr\rangle+\sum_p\mathrm c_p\mathfrak Q_p\bigl(\mathring\omega^{(n)}\bigr)+\sum_{a\in\mathfrak A_n}\beta_a\E\bigl[E_{a,o}\bigl(k;\mathring\omega^{(n)}\bigr)\bigr]\Bigr]\\
&+\mathbf m\sum_{n=0}^{n_1-1}\mathfrak U\bigl(\mathring\omega^{(n+1)},\mathring\omega^{(n)}\bigr),
\end{aligned}
\label{eq:link-history-objective}
\end{equation}
which, when the agents' latent law does not read the links, is the generator part of the expected
negative log joint density. Its minimizers are the most
probable generator histories given the matter recognition, in the sense of
\cref{prop:link-map-inertia}. For one episode with every $\beta_a=1$, and an agent latent law that
does not read the links, $\mathscr S^{\mathrm{hist}}$ is the function $\mathbf V$ of
\cref{prop:link-field-equation}(d) up to a constant, so $\nabla\mathscr S^{\mathrm{hist}}=0$ is the history
form of the most-probable-link condition $\nabla\mathbf V=0$. On $\mathscr T^{n_1+1}$, with
$\langle\mathring\omega,\mathring\omega'\rangle=\sum_n\langle\mathring\omega^{(n)},\mathring\omega'^{(n)}\rangle$,
define self-adjoint operators $\mathbf K$ and $\mathbf N_0$ and a source $\mathbf j$ by
\begin{equation}
\begin{aligned}
\langle\mathring\omega,\mathbf K\mathring\omega\rangle
={}&\sum_n\Bigl[\bigl\langle\mathring\omega^{(n)},\mathring\Sigma^{-1}\mathring\omega^{(n)}\bigr\rangle+\sum_p\mathrm c_p\bigl\lVert\mathbf C_p\mathring\omega^{(n)}\bigr\rVert_{\mathsf G}^2
+\sum_{a\in\mathfrak A_n}\beta_a\Tr\bigl(W_a\mathrm d\rho\mathring\omega^{(n)}_e\E\bigl[k_jk_j^\top\bigr](\mathrm d\rho\mathring\omega^{(n)}_e)^\top\bigr)\Bigr]\\
&+\mathbf m\sum_{n<n_1}\bigl\lVert\mathring\omega^{(n+1)}-\mathring\omega^{(n)}\bigr\rVert^2,\\
\langle\mathring\omega,\mathbf N_0\mathring\omega\rangle
={}&-\sum_n\sum_{a\in\mathfrak A_n}\beta_a\Tr\bigl(\mathbf J_a^\top(\mathrm d\rho\mathring\omega^{(n)}_e)^2\bigr),
\qquad
\langle\mathbf j,\mathring\omega\rangle=\sum_n\sum_{a\in\mathfrak A_n}\beta_a\Tr\bigl(\mathbf J_a^\top\mathrm d\rho\mathring\omega^{(n)}_e\bigr),
\end{aligned}
\label{eq:link-linear-operator}
\end{equation}
where $e$ and $j$ are those of $a$, $\E[k_jk_j^\top]=\Sigma^{(n)}_j+\mu^{(n)}_j\mu^{(n)\top}_j$,
$\mathbf C_p\mathring\omega=(\mathbf C\mathring\omega)_p$ is the loop sum of the cycle map
\eqref{eq:link-cycle-map}, and, for $a\in\mathfrak A_n$,
\begin{equation}
\mathbf J_a=W_a\bigl[(\mu^{(n)}_i-\mu^{(n)}_j)\mu^{(n)\top}_j-\Sigma^{(n)}_j\bigr]
\label{eq:link-zero-current}
\end{equation}
is the current \eqref{eq:link-delta-current} of record $a$ at zero generators for a product law. Thus
$\mathbf K=\mathring\Sigma^{-1}+\sum_p\mathrm c_p\mathbf C_p^\top\mathbf C_p+\mathbf m(-\Delta_n)+\mathring\Phi$,
where $\mathring\Sigma^{-1}$ and the loop term act episode by episode,
$(-\Delta_n\mathring\omega)^{(n)}=2\mathring\omega^{(n)}-\mathring\omega^{(n+1)}-\mathring\omega^{(n-1)}$ in the
interior, $\mathring\omega^{(0)}-\mathring\omega^{(1)}$ and $\mathring\omega^{(n_1)}-\mathring\omega^{(n_1-1)}$ at
the ends, and $\mathring\Phi$ is the matter operator of the last term of the bracket. Let
$\mathrm c_{\mathring\Sigma}$ be the smallest eigenvalue of $\mathring\Sigma^{-1}$, and let
$\operatorname{cond}(\mathsf G^{\mathrm{fib}})$ be the ratio of the extreme eigenvalues of
$\mathsf G^{\mathrm{fib}}$.
\begin{enumerate}
\item[(a)] $\nabla\mathscr S^{\mathrm{hist}}(0)=-\mathbf j$ and $\nabla^2\mathscr S^{\mathrm{hist}}(0)=\mathbf K+\mathbf N_0$, so to first order in
the generators and the currents the stationarity condition is
\begin{equation}
\mathbf K\mathring\omega=\mathbf j.
\label{eq:link-linearized-equation}
\end{equation}
\item[(b)] $\mathbf K\succeq\mathrm c_{\mathring\Sigma}I$. Hence \eqref{eq:link-linearized-equation}
has the unique solution $\mathring\omega_{\mathrm{lin}}=\mathbf K^{-1}\mathbf j$, with
$\lVert\mathring\omega_{\mathrm{lin}}\rVert\leq\lVert\mathbf j\rVert/\mathrm c_{\mathring\Sigma}$, and
$\lVert\mathbf K^{-1}\mathbf j-\mathbf K^{-1}\mathbf j'\rVert\leq\lVert\mathbf j-\mathbf j'\rVert/\mathrm c_{\mathring\Sigma}$.
\item[(c)] Let $\mathrm c_0=\lVert\mathbf N_0\rVert$ and, for $\mathrm r>0$,
$\mathrm c_{\mathrm r}=\sup_{\lVert\mathring\omega\rVert\leq\mathrm r}\lVert D^3\mathscr S^{\mathrm{hist}}(\mathring\omega)\rVert$. Then
$\mathrm c_0\leq\operatorname{cond}(\mathsf G^{\mathrm{fib}})\max_{n,e\in E_+}\mathrm j^{(n)}_e$ and
$\lVert\mathbf j\rVert\leq\operatorname{cond}(\mathsf G^{\mathrm{fib}})^{1/2}(\sum_{n,e\in E_+}(\mathrm j^{(n)}_e)^2)^{1/2}$,
where $\mathrm j^{(n)}_e$ is the sum of $\beta_a\lVert\mathbf J_a\rVert_{\mathrm F}$ over the records
$a\in\mathfrak A_n$ on $e$ or $\bar e$. If $\mathrm c_0+\mathrm c_{\mathrm r}\mathrm r\leq\mathrm c_{\mathring\Sigma}/2$ and
$\lVert\mathbf j\rVert\leq\mathrm c_{\mathring\Sigma}\mathrm r/4$, then $\mathscr S^{\mathrm{hist}}$ has exactly one critical point
$\mathring\omega_\star$ in the ball of radius $\mathrm r$ about $0$; it is the unique minimizer of $\mathscr S^{\mathrm{hist}}$ on
that ball, $\lVert\mathring\omega_\star\rVert\leq2\lVert\mathbf j\rVert/\mathrm c_{\mathring\Sigma}$, and
\begin{equation}
\lVert\mathring\omega_\star-\mathring\omega_{\mathrm{lin}}\rVert
\leq\frac{2\lVert\mathbf j\rVert}{\mathrm c_{\mathring\Sigma}^2}\Bigl(\mathrm c_0+\frac{\mathrm c_{\mathrm r}\lVert\mathbf j\rVert}{\mathrm c_{\mathring\Sigma}}\Bigr),
\label{eq:link-linearization-error}
\end{equation}
which, at fixed $\mathrm c_{\mathrm r}$, is of second order in the currents.
\end{enumerate}
\status{ESTABLISHED}

The proof is in \cref{sec:proofs-links}.

In the language of lattice gauge theory, \cref{prop:link-linearized-equation} reads as follows, and
the reading adds nothing to it. The spatial term $\sum_p\mathrm c_p\mathbf C_p^\top\mathbf C_p$ is a
discrete curl-curl operator: $\mathbf C_p$ sums the generators around the loop $p$, so the term
penalizes the circulation of the generator field on the cycle space of the multigraph, the magnetic
part of a lattice action. The episode term $\mathbf m(-\Delta_n)$ has the quadratic form
$\mathbf m\sum_n\lVert\mathring\omega^{(n+1)}-\mathring\omega^{(n)}\rVert^2$ of a forward difference in the
episode index; by \cref{prop:link-history-prior}(c) and \cref{prop:link-map-inertia}(b) it is the
second-order part of the time-like loop defect with identity time-like links, the temporal gauge, in
which the time-like plaquettes give the electric term of the Hamiltonian of
\citet{KogutSusskind1975}, as derived in \citet[Ch.~15, Eq.~(15.41)]{Creutz1983}. The anchor
precision $\mathring\Sigma^{-1}$ is a bare mass and the matter operator $\mathring\Phi$ a state-dependent
screening mass, so at linear order $\mathbf K$ has the form of a massive Euclidean vector-field
operator of Proca type, a mass term plus curl-curl plus a squared episode difference, acting on the
spatial components alone, since the history carries no time-like variable. Its curl-curl and episode
terms act as scalar graph operators on each component of $\mathfrak g$ in an orthonormal basis, and
the bracket of $\mathfrak g$ enters only through the terms of order three and higher, which
\cref{prop:link-linearized-equation}(c) bounds. The episode index is an
inference index in the sense of \cref{sec:intro-conventions}, and the signature is Euclidean
because $\mathbf K$ is positive (\cref{prop:link-linearized-equation}(b)).

The loop term sees only part of the link space, and the declared family $\mathbb L$ fixes which part.
The next statement makes this exact for the quadratic loop action of
\cref{prop:link-anchor-m-coordinate}(c) and for the loop term of $\mathbf K$.

\corollaryheading{Cycle-space splitting of the loop sector}{cor:link-cycle-splitting}
Let $\mathfrak I=(V,E,\bar{\ })$ be the interaction multigraph of \cref{def:geo-graph-links}, with
$\kappa$ connected components, write $\mathscr T=\R^{E_+}\otimes\mathfrak g$, let
$\mathbb L_+=\{p\in\mathbb L:\mathrm c_p>0\}$, and let $\mathbf C_+:\mathscr T\to\mathfrak g^{\mathbb L_+}$ be the
cycle map \eqref{eq:link-cycle-map} of $\mathbb L_+$. Let $d:\R^V\to\R^{E_+}$ be $(da)_e=a_i-a_j$ for
$e:j\to i$, let $Z=(\operatorname{im}d)^\perp$ be the cycle space, and for $p\in\mathbb L_+$ let
$z_p\in\R^{E_+}$ count the traversals of each $e\in E_+$ along its orientation minus those against it,
so that $(\mathbf C\mathring\omega)_p=\sum_{e\in E_+}z_{p,e}\mathring\omega_e$. Put
$r_{\mathbb L}=\dim\operatorname{span}\{z_p:p\in\mathbb L_+\}\leq\lvert\mathbb L_+\rvert$, and call
\begin{equation}
\mathscr T_{\mathrm{ex}}=\operatorname{im}d\otimes\mathfrak g,
\qquad
\mathscr T_{\mathrm{har}}=\bigl(Z\cap\{z_p:p\in\mathbb L_+\}^\perp\bigr)\otimes\mathfrak g
\label{eq:link-exact-harmonic}
\end{equation}
the exact and the harmonic generators; $\mathring\omega$ is exact when $\mathring\omega_e=a_i-a_j$ for some
$a\in\mathfrak g^V$.
\begin{enumerate}
\item[(a)] $\ker\mathbf C_+=\mathscr T_{\mathrm{ex}}\oplus\mathscr T_{\mathrm{har}}$, a sum orthogonal for the
inner product of $\mathscr T$, with
\begin{equation}
\dim\mathscr T_{\mathrm{ex}}=\dim\mathfrak g\,(\lvert V\rvert-\kappa),
\qquad
\dim\mathscr T_{\mathrm{har}}=\dim\mathfrak g\,(\lvert E_+\rvert-\lvert V\rvert+\kappa-r_{\mathbb L})
\geq\dim\mathfrak g\,(\lvert E_+\rvert-\lvert V\rvert+\kappa-\lvert\mathbb L_+\rvert).
\label{eq:link-cycle-splitting-dims}
\end{equation}
In the translation instance (\cref{prop:link-translation-instance}) an exact generator shift by
$a_i-a_j$ is the link change $\Theta_e\mapsto\Theta_e+a_i-a_j$ with the frames held, and every loop
defect $\mathfrak Q_p$ with $p\in\mathbb L_+$ is constant along $\ker\mathbf C_+$.
\item[(b)] Under the quadratic action of \cref{prop:link-anchor-m-coordinate}(c), with
$\mathring\Sigma=\bigoplus_{e\in E_+}\mathring\Sigma_e$,
\begin{equation}
\mathring\alpha^{\mathrm{new}}-\E_R[\mathring\omega]\in\mathring\Sigma\operatorname{im}\mathbf C_+^\top,
\qquad
\mathscr T=\mathscr T_{\mathrm{ex}}\oplus\mathscr T_{\mathrm{har}}\oplus\mathring\Sigma\operatorname{im}\mathbf C_+^\top,
\label{eq:link-anchor-coexact}
\end{equation}
and the last summand is orthogonal to the first two for
$\langle\cdot,\mathring\Sigma^{-1}\cdot\rangle$. Hence $\mathring\alpha^{\mathrm{new}}$ and $\E_R[\mathring\omega]$ have
the same components in $\mathscr T_{\mathrm{ex}}$ and in $\mathscr T_{\mathrm{har}}$, and the anchor
update corrects only the component in $\mathring\Sigma\operatorname{im}\mathbf C_+^\top$.
\item[(c)] For the operator $\mathbf K$ of \cref{prop:link-linearized-equation}, if every
$\mathring\omega^{(n)}$ lies in $\ker\mathbf C_+$, the loop term vanishes and
$\langle\mathring\omega,\mathbf K\mathring\omega\rangle=\sum_n\langle\mathring\omega^{(n)},\mathring\Sigma^{-1}\mathring\omega^{(n)}\rangle
+\mathbf m\sum_{n<n_1}\lVert\mathring\omega^{(n+1)}-\mathring\omega^{(n)}\rVert^2+\langle\mathring\omega,\mathring\Phi\mathring\omega\rangle$.
\end{enumerate}
\status{ESTABLISHED}

\noindent\emph{Proof.}
(a) The kernel of $d$ is the functions constant on components, so
$\dim\operatorname{im}d=\lvert V\rvert-\kappa$ and $\dim Z=\lvert E_+\rvert-\lvert V\rvert+\kappa$. Along a
closed walk every arrival at a vertex is followed by a departure, so the net traversals at each vertex
cancel, $d^\top z_p=0$, and $z_p\in Z$. For the matrix $\mathsf Z_+$ with rows $z_p$, $p\in\mathbb L_+$,
this gives $\ker\mathsf Z_+=(\operatorname{span}\{z_p\})^\perp=\operatorname{im}d\oplus\bigl(Z\cap\{z_p\}^\perp\bigr)$,
orthogonally, and the second summand has dimension $\dim Z-r_{\mathbb L}$ because
$\operatorname{span}\{z_p\}\subset Z$. By \eqref{eq:link-cycle-map} and the reversal convention
\eqref{eq:link-generator-reversal}, $\mathbf C_+=\mathsf Z_+\otimes\mathrm{id}_{\mathfrak g}$, and the inner product
of $\mathscr T$ is the Euclidean one on $\R^{E_+}$ times $\langle\cdot,\cdot\rangle_{\mathsf G}$, so tensoring with
$\mathfrak g$ carries the decomposition, its orthogonality, and the dimensions over. In the translation
instance $\mathring\omega_e=\Theta_e-U_i+U_j$ by \cref{prop:link-translation-instance}(a), so with the frames
held an exact shift of $\mathring\omega$ is the stated shift of $\Theta$, and by
\eqref{eq:link-translation-defect} $\mathfrak Q_p=\tfrac12\lVert(\mathbf C\mathring\omega)_p\rVert_{\mathsf G}^2$
depends on $\mathring\omega$ only through $\mathbf C_+\mathring\omega$ when $p\in\mathbb L_+$.
(b) Terms with $\mathrm c_p=0$ vanish in \eqref{eq:link-anchor-quadratic}, so
$\mathring\alpha^{\mathrm{new}}-\E_R[\mathring\omega]=\mathring\Sigma\sum_{p\in\mathbb L_+}\mathrm c_p\mathbf C_p^\top\mathbf C_p\E_R[\mathring\omega]\in\mathring\Sigma\operatorname{im}\mathbf C_+^\top$.
Each $\mathring\Sigma_e$ is self-adjoint and positive, so for $h\in\ker\mathbf C_+$ and
$v\in\operatorname{im}\mathbf C_+^\top$, $\langle\mathring\Sigma v,\mathring\Sigma^{-1}h\rangle=\langle v,h\rangle=0$.
The subspace $\mathring\Sigma\operatorname{im}\mathbf C_+^\top$ therefore meets $\ker\mathbf C_+$ only in zero, and its
dimension $\operatorname{rank}\mathbf C_+=\dim\mathscr T-\dim\ker\mathbf C_+$ makes the sum all of $\mathscr T$; with
(a) this gives \eqref{eq:link-anchor-coexact}, and the components of the difference in the first two
summands vanish. (c) $\mathbf C_p\mathring\omega^{(n)}=0$ for $p\in\mathbb L_+$, and the other loop terms carry
$\mathrm c_p=0$, so only the remaining terms of \eqref{eq:link-linear-operator} survive. $\square$

Attaching a two-cell along each walk of $\mathbb L_+$ makes the declared loop family a two-complex on
the agents, and \cref{cor:link-cycle-splitting} is the Hodge decomposition of its one-cochains with
values in $\mathfrak g$ \citep[Theorem~3.1]{HansenGhrist2019}, applied to the constant sheaf: exact
generators are its gradient part, $\mathring\Sigma\operatorname{im}\mathbf C_+^\top$ its circulating part in the
metric of the prior, and $\mathscr T_{\mathrm{har}}$ the circulation around cycles of $\mathfrak I$ that bound
no declared walk. By \eqref{eq:link-cycle-splitting-dims}, a family of short loops on a multigraph whose
cycle space is mostly spanned by longer cycles leaves most of that space to the anchor, the episodes,
and the matter, so the choice of $\mathbb L$ is a structural declaration on the agent set of the same
kind as the support of the source prior. Away from zero generators the loop defects of a nonabelian
group are not functions of $\mathbf C\mathring\omega$, and the splitting describes only their quadratic part
(\cref{prop:link-loop-defect}(e)).

\section*{Physicist's summary}

The links between agents are treated here as a field that is inferred, not as a background that is held.
Each link is the exponential of a Lie-algebra generator measured in the frame of a reference point;
per-agent frame changes do not touch the generators, and a change of reference point rotates all of them
by one adjoint map. A Gaussian law on the generators, reweighted by a Wilson-type loop penalty, is a
normalized prior for any structure group, so no finite Haar measure is needed; on the general linear
group the loop penalty is measured with a fiber metric, because conjugation-invariant penalties cannot
see unipotent holonomy. The links enter the free energy through one relative entropy from this prior.
Stationarity is a field equation: the restoring force of the prior balances a matter current, which for
a comparison record is a precision-weighted prediction error times the source state, the form of a delta
rule, transported by the derivative of the exponential map. Shrinking the prior to its anchor recovers
the free energy and the evidence of the held links at the anchor. Integrating out matter that is normalized jointly with the links generates an induced
link action, as in induced lattice gauge theory, while matter attached by record extension generates
none. Inference of the generators is first-order gradient flow in the Fisher geometry of their Gaussian
law, with the Bayesian learning rule as its Euler step; inertia appears only in the most probable
history across learning episodes, where the precision of the increment acts as a mass. Linearized about
zero generators, the stationarity operator is a bare mass from the anchor, a curl-curl term from the
loops, a squared episode difference, and a screening mass from the matter, which is the form of a
massive Euclidean vector operator in the episode index. On the cycle space the loop term acts only on the circulating part of the
generators: exact generators and circulation around cycles that bound no declared loop are fixed by the
anchor, the episodes, and the matter alone, and the anchor update corrects only the circulating part.
